% v4.8.0 post-release layers
% Appended layers v4.8.0-r03 to r20, carried verbatim from the full-monolith working source
% Quantum_Sphaera_Companion_v4_8_0_source.tex (SHA256 a1b988d3c9e3642f93e201bcb56d6ff03020ae6f284339f29c7030f7b092ebe5),
% lines 65055-67717.  Not a preserved v4.7 fragment; earlier layers are not rewritten.
% Content SHA256 of the carried block: cc1ee73f98622f6aa7bb00857507a3b316b1bbb754111dd6b4728930ec4d6754

% ============================================================================
\clearpage
\section*{v4.8.0-r03 Source-Only: Dimensional and Numerical Audit of the Physical-Constant Layer}
\addcontentsline{toc}{section}{v4.8.0-r03 Source-Only: Dimensional and Numerical Audit of the Physical-Constant Layer}
\label{sec:v480-r03-dimensional-audit}
% ============================================================================

\paragraph{Purpose.}
An external structural evaluation of v4.7.0 rated equation and dimensional integrity as the
weakest of its five criteria.  This layer audits the places where the monolith assigns
physical units, CODATA values or frame-insertion counts: the nomenclature, the
phase-shifted constants of \S\ref{sec:phase-constants}, the Hawking analysis of
\S\ref{sec:bh}, the cosmological statements around
Theorems~\ref{thm:dark-offset} and~\ref{thm:dark-energy}, and the photon-frame
\(\Lambda\) block of \S\ref{sec:lambda_cancel}.  Every number was recomputed.  Earlier layers
are not rewritten; each finding is recorded here with the correction that follows from the
document's own rules, or with the reason why no such correction exists.

\noindent\textbf{Source-only status.}
This v4.8.0-r03 layer is source-only.  It does not build a PDF, does not assign a new DOI,
and does not authorize a public release.  The v4.8.0 Version DOI
\href{https://doi.org/10.5281/zenodo.23365244}{10.5281/zenodo.23365244} and the Concept DOI
\href{https://doi.org/10.5281/zenodo.19210728}{10.5281/zenodo.19210728} are unchanged by it.

\subsection*{v4.8.0-r03 the frame-insertion calculus}
\label{subsec:v480-r03-k-calculus}

The definition in \S\ref{sec:phase-constants} assigns to a constant \(X\) the photon-frame
value \(X^{(\mathrm{photon})}=\gL^{-k(X)}X^{(\mathrm{matter})}\).

\paragraph{Lemma v4.8.0-r03.1.}
If \(X=\prod_i Y_i^{a_i}\) with real exponents \(a_i\) and every factor obeys the rule, then
\(X\) obeys it with
\[
k(X)=\sum_i a_i\,k(Y_i).
\]

\paragraph{Proof.}
\(X^{(\mathrm{photon})}=\prod_i\bigl(\gL^{-k(Y_i)}Y_i^{(\mathrm{matter})}\bigr)^{a_i}
=\gL^{-\sum_i a_ik(Y_i)}X^{(\mathrm{matter})}\).  \(\square\)

The lemma is what Theorem~\ref{thm:planck-photon} uses: \(\lP=(\hbar G_N/c^3)^{1/2}\) gives
\(k(\lP)=\tfrac12(1+2-0)=\tfrac32\), as Proposition~\ref{prop:uv-photon} also states.  Since
the rule multiplies dimensionful quantities by the pure number \(\gL\), it is dimensionally
consistent by construction.  What it requires is that every composite quantity be assigned
the \(k\) the lemma gives, and that every factor of the composite carry its own \(k\).  The
findings below are the places where that is not done.

\subsection*{v4.8.0-r03 findings}
\label{subsec:v480-r03-findings}

\paragraph{F1 --- vacuum-energy ratio (\(k\)-calculus).}
Equation~\eqref{eq:rho-ratio}, Corollary~\ref{cor:matter-error} and the corollary
``Cosmological constant revisited'' in \S\ref{sec:phase-constants} use
\(\rho_{\mathrm{vac}}\sim\lP^{-4}\) and obtain the frame ratio \(\gL^{6}\).  The relation
\(\rho\sim\lP^{-4}\) holds only in units with \(\hbar=c=1\).  The framework assigns
\(k(\hbar)=1\), so \(\hbar\) cannot be set to one in both frames.  As an energy density,
\(\rho_{\mathrm{vac}}\sim\hbar c\,\lP^{-4}\), and Lemma~v4.8.0-r03.1 gives
\[
k(\rho_{\mathrm{vac}})=1+0-4\cdot\tfrac32=-5,
\qquad
\frac{\rho_{\mathrm{vac}}^{(\mathrm{photon})}}{\rho_{\mathrm{vac}}^{(\mathrm{matter})}}=\gL^{5}=22.53\ldots,
\]
not \(\gL^{6}\).  The mass-density form \(m_P/\lP^{3}\) gives the same \(k=-\tfrac12-\tfrac92=-5\).
In addition, the source quotes \(\gL^{6}\) as \(41.8\) in one place and \(42.0\) in another;
with \(\gL=1.86444\ldots\) the value is \(42.004\ldots\).

\paragraph{F2 --- fine-structure drift (sign of the exponent and empirical status).}
With \(k(\alpha^{-1})=1\) (Theorem~\ref{thm:k-values}) the rule gives
\((\alpha^{-1})^{(\mathrm{photon})}=\gL^{-1}(\alpha^{-1})^{(\mathrm{matter})}\), hence
\[
\alpha^{(\mathrm{matter})}=\gL^{-1}\,\alpha^{(\mathrm{photon})},
\]
whereas the corollary ``\(\alpha\)-drift as a frame effect'' writes
\(\alpha^{(\mathrm{matter})}(t)=\gL^{(\mathrm{local})}(t)\,\alpha^{(\mathrm{photon})}\).  The two
are inverse to each other.  Separately, the corollary speaks of an ``observed NIST drift''.  The
NIST Al\(^+\)/Hg\(^+\) clock comparison of 2008 found a drift consistent with zero, and a fit of
the same data for an annual modulation gives a fractional amplitude of about
\((-8.9\pm8.4)\times10^{-18}\), also consistent with zero.  There is therefore no observed drift
to explain.  The annual-modulation prediction survives only as a constraint:
\(|\Delta\gL^{(\mathrm{local})}/\gL|\lesssim10^{-17}\) over the year.

\paragraph{F3 --- three incompatible types for \(\Lambda\).}
The monolith uses \(\Lambda\) in three dimensionally different roles:
\[
\begin{array}{ll}
\Lambda_{\mathrm{bare}}\sim\gamma_1^2\approx200\,m_P^4 & (\text{Prop.~\ref{prop:bare-lambda}}),\\
\Lambda_{\mathrm{eff}}\sim\bsym^2\rho_P & (\text{near Thm.~\ref{thm:dark-offset}}),\\
\Lambda_{\mathrm{obs}}=\bsym^2\,(\lP^{(\mathrm{photon})})^{-2} & (\text{\S\ref{sec:falsifiable}}).
\end{array}
\]
The first equates a pure number with a Planck-unit energy density, the second is an energy
density, the third a curvature with unit \(\mathrm{m}^{-2}\).  Curvature and energy density are
related by \(\Lambda=8\pi G_N\rho_\Lambda/c^4\); the factor \(8\pi\) is absent throughout.  The
photon-frame step of Proposition~\ref{prop:bare-lambda} uses \(k(\Lambda)=2\), while the
curvature form \(\Lambda\propto\lP^{-2}\) gives \(k=-3\) under Lemma~v4.8.0-r03.1.  The source
already states that a self-consistent \(k(\Lambda)\) is open; r03 records that the value \(2\)
used in the display is not the one the calculus gives for any of the three forms.

\paragraph{F4 --- the ``unit-system artifact'' remark.}
Remark~\ref{rem:open} and the prediction P4 of \S\ref{sec:falsifiable} state that the
\(10^{-122}\) of the cosmological-constant problem is an artefact of SI units, and that in
natural units \(\Lambda_{\mathrm{obs}}\sim\bsym^2\approx0.30\).  This is a dimensional error.
The quantity
\[
\Lambda_{\mathrm{obs}}\,\lP^{2}=1.1056\times10^{-52}\,\mathrm{m}^{-2}\cdot(1.6163\times10^{-35}\,\mathrm{m})^{2}
=2.89\times10^{-122}
\]
is dimensionless and takes the same value in every unit system; in units with \(\lP=1\) the
measured \(\Lambda_{\mathrm{obs}}\) is \(2.89\times10^{-122}\), not \(0.30\).  Taken literally,
\(\Lambda_{\mathrm{obs}}=\bsym^2(\lP^{(\mathrm{photon})})^{-2}=7.5\times10^{69}\,\mathrm{m}^{-2}\)
exceeds the measured value by a factor \(6.8\times10^{121}\).  P4 is therefore contradicted by
measurement as stated.  The same block also gives three different sizes for the required
cancellation (\(2.3\times10^{-4}\), \(1.5\times10^{-3}\) and a factor \(3.1\)), which cannot all
describe one ratio.

\paragraph{F5 --- cosmological fractions.}
Theorem~\ref{thm:dark-energy} writes \(1-\bsym^2\approx\tfrac23\approx0.698\); with
\(\bsym^2=0.30178\ldots\) the value is \(0.69822\ldots\), which is not \(\tfrac23\).  Two different
predictions for \(\Omega_\Lambda/\Omega_M\) occur: \(\bsym^2/\gstring=2.37\) near
Theorem~\ref{thm:dark-offset} and \((1-\bsym^2)/\bsym^2=2.31\) from
Theorem~\ref{thm:dark-energy}.  They are different functions of the same parameters and cannot
both hold.  The observed ratio from the Planck 2018 values quoted in
Remark~\ref{rem:omega-bsym} is \(0.6889/0.3111=2.214\), not \(2.18\).  The same remark writes
\(\bsym=\gBI/\gstring\approx0.5493\); that quotient is \(\gL=1.864\), and the stated value
belongs to \(\bsym=\sqrt{(\gL-1)/(\gL+1)}\).

\paragraph{F6 --- Hawking temperature.}
Proposition~\ref{prop:hawking-gBI-tachyon} derives
\(T_H=\hbar c^3/(8\pi G_NMk_B)=m_P^2c^2/(8\pi Mk_B)\) correctly.  The step to
\(k(T_H)=-1\) declares \(k(k_B)=0\) but also needs \(k(M)=0\) for the black-hole mass, which is
not declared, and which differs from the assignment \(k(m_ec^2)=1\) for the electron.  Under
Lemma~v4.8.0-r03.1, \(k(T_H)=1-2-k(M)=-1-k(M)\).

\paragraph{F7 --- notation.}
In Theorem~\ref{thm:dark-offset} the letter \(k\) denotes a spring constant, while \(k(X)\)
denotes the frame-insertion count elsewhere in the same part.  The oscillator frequency
\(\omega=\gstring\sqrt{2/3}\) is a pure number; Planck units are implied but not declared.

\paragraph{Verified without finding.}
\(\gstring=\ln2/(\pi\sqrt3)=0.12738\ldots\);
\(\gL=\gBI/\gstring=1.86444\ldots\);
\(\bsym=\sqrt{(\gL-1)/(\gL+1)}=0.54935\ldots\);
\(\gL^{3/2}=2.5458\ldots\);
\(\lP^{(\mathrm{photon})}=6.35\times10^{-36}\,\mathrm{m}\);
\(k(\lP)=\tfrac32\);
the units of \(G_N\), \(\hbar\), \(c\) in Theorem~\ref{thm:k-values};
\(T_H=\hbar c^3/(8\pi G_NMk_B)\);
\((\alpha^{-1})^{(\mathrm{photon})}=73.50\ldots\) under \(k(\alpha^{-1})=1\).

\paragraph{Theoremlet v4.8.0-r03.}
Under the frame-insertion rule of \S\ref{sec:phase-constants} and Lemma~v4.8.0-r03.1:
\begin{enumerate}[leftmargin=2em]
 \item the photon-to-matter ratio of a Planck-cutoff vacuum energy density is \(\gL^{5}\), not
 \(\gL^{6}\);
 \item \(k(\alpha^{-1})=1\) implies \(\alpha^{(\mathrm{matter})}=\gL^{-1}\alpha^{(\mathrm{photon})}\);
 \item \(\Lambda\) has no single \(k\) until one of its three forms is chosen, and none of them
 gives \(k=2\);
 \item no choice of units turns the dimensionless number \(\Lambda_{\mathrm{obs}}\lP^2\approx
 2.9\times10^{-122}\) into a number of order one.
\end{enumerate}

\paragraph{Proof sketch.}
Items 1--3 are applications of the lemma to \(\hbar c\,\lP^{-4}\), to \(\alpha=(\alpha^{-1})^{-1}\),
and to \(\lP^{-2}\), \(\hbar c\,\lP^{-4}\) and \(G_N\rho/c^4\) respectively.  Item 4 holds because a
ratio of two quantities of the same dimension is invariant under change of units.  \(\square\)

\subsection*{v4.8.0-r03 re-check of P4}
\label{subsec:v480-r03-p4-recheck}

F4 was checked a second time, against the main paper and against a reading of P4 that is
dimensionally consistent.

\paragraph{Origin of the unit \(\lP^{-2}\).}
The main paper states the frame contribution without a unit.  Its Theorem~11.6 (Dark energy as
permanent frame offset) gives \(\Lambda_{\mathrm{eff}}\sim\bsym^2\rho_P\), its Proposition~12.4
(Frame term gives the cosmological constant) gives \(\Lambda_{\mathrm{obs}}\sim\bsym^2\), and its
remark on the cosmological-constant problem records that reducing the bare value by
\(10^{122}\) through the \(\Theta\)--\(R\) cancellation ``is not yet established''
\cite{brendecke2026main}.  The unit \(\lP^{-2}\) and the statement that \(10^{-122}\) is a unit
artefact enter only in the photon-frame \(\Lambda\) block of the companion,
\S\ref{sec:lambda_cancel}.  The error recorded in F4 therefore lies in that later
reinterpretation, not in the main paper.

\paragraph{F8 --- ``dimensionally correct''.}
The proof of Theorem~\ref{thm:lambda-zero}(iii) calls \(\Lambda_{\mathrm{obs}}=\bsym^2/\lP^{(\mathrm{matter})2}\)
``dimensionally correct up to the \(\gL^{3/2}\) Planck-length correction''.  The unit
\(\mathrm{m}^{-2}\) is correct; the value is \(1.2\times10^{69}\,\mathrm{m}^{-2}\), a factor
\(10^{121}\) above the measurement, and no power of \(\gL\) closes that gap.  The theorem uses the
matter-frame Planck length, P4 the photon-frame one.

\paragraph{F9 --- citation.}
The proof of Theorem~\ref{thm:lambda-zero}(iii) and the status table of the \(\Lambda\) block cite
``Proposition~18.4 of the main paper''.  The main paper has no numbered environment 18.4.  The
statement meant is its Proposition~12.4, which asserts only the proportionality
\(\Lambda_{\mathrm{obs}}\sim\bsym^2\).

\paragraph{F10 --- units of the balance condition.}
In Equation~\eqref{eq:tr-balance} the integrand \(\Sigma(\tau)\) contains \(e^{-\tau/2}\), so
\(\tau\) is a pure number, while the upper limit \(\lP^{(\mathrm{photon})}\) and the right-hand
side \(\bsym^2\lP^{(\mathrm{photon})2}\) are written as a length and a squared length.  The
condition is consistent only when read in matter-frame Planck units, where
\(\lP^{(\mathrm{photon})}=\gL^{-3/2}=0.3928\ldots\) and the right-hand side is
\(\bsym^2\gL^{-3}=0.0466\ldots\).  This reading is not stated in the source.

\paragraph{A dimensionally consistent reading, offered and not adopted.}
Since \(\bsym^2\) is a pure number, the consistent place for it is a dimensionless fraction of
the critical density, not \(\Lambda\) itself.  If the Markov statement of
Theorem~\ref{thm:dark-energy}, \(\Omega_\Lambda=1-\bsym^2\), is taken as the prediction, then
\[
\Lambda_{\mathrm{obs}}=\frac{3\,\Omega_\Lambda H_0^2}{c^2}
=\frac{3\,(1-\bsym^2)\,H_0^2}{c^2}
=1.112\times10^{-52}\,\mathrm{m}^{-2}
\qquad(H_0=67.4\ \mathrm{km\,s^{-1}\,Mpc^{-1}}),
\]
which differs from \(\Lambda_{\mathrm{obs}}=1.1056\times10^{-52}\,\mathrm{m}^{-2}\) by \(0.6\%\).
At the level of the fraction, \(1-\bsym^2=0.6982\ldots\) compares with the Planck 2018 values
\(0.6847\) and \(0.6889\) by \(2.0\%\) and \(1.4\%\).

The reading has a price.  It predicts the fraction \(\Omega_\Lambda\), and takes the absolute scale
from the measured \(H_0\).  The number \(10^{-122}\) does not disappear: it becomes
\(\Lambda_{\mathrm{obs}}\lP^2=3\Omega_\Lambda(\lP/R_H)^2\) with the Hubble radius \(R_H=c/H_0\), so
that the remaining question is the hierarchy \(R_H/\lP=8.49\times10^{60}\).  The source
quantity closest to it is the active zero count \(n_H\sim10^{62}\) at the Hubble time
(Theorem~11.5 of the main paper); it differs from \(R_H/\lP\) by a factor of about \(12\), and the
source does not identify the two.  The reading also selects \(1-\bsym^2\) among the two
conflicting predictions of F5.

Whether P4 is restated in this form, restated otherwise, or withdrawn is left to the author.
r03 records the reading as an option and asserts nothing beyond the arithmetic above.

\subsection*{v4.8.0-r03 provenance of the frame parameters}
\label{subsec:v480-r03-frame-parameter-provenance}

\paragraph{F11 --- both inputs of \(\gL\) are values of the same LQG parameter.}
The frame factor \(\gL=\gBI/\gstring=1.864\ldots\) enters every frame conversion of the monolith,
\(\bsym=\sqrt{(\gL-1)/(\gL+1)}\), the cosmological fractions of
Theorem~\ref{thm:dark-energy}, the photon-frame Planck length of
Theorem~\ref{thm:planck-photon} and the instance \(\lambda_\gamma=\gL^{-3/2}\) of
v4.8.0-r01.  Its two inputs have the following origin in the literature.
\begin{enumerate}[leftmargin=2em]
 \item \(\gstring=\ln2/(\pi\sqrt3)=0.12738\ldots\), introduced in the nomenclature as
 ``string coupling'', is the value of the Barbero--Immirzi parameter obtained by Ashtekar, Baez,
 Corichi and Krasnov from the state counting of quantum isolated horizons
 (arXiv \texttt{gr-qc/9710007}, \texttt{gr-qc/0005126}).  It is not a string-theory coupling.
 \item \(\gBI\approx0.2375\) is the value of the same parameter after the corrected state
 counting of Domagala and Lewandowski and of Meissner (arXiv \texttt{gr-qc/0407051},
 \texttt{gr-qc/0407052}).  Domagala and Lewandowski show that the earlier counting contains a
 spurious constraint that suppresses spins above \(j=\tfrac12\), and that the earlier value
 therefore does not hold.
\end{enumerate}
Hence \(\gL\) is the ratio of the corrected to the superseded determination of one parameter.
Theorem~\ref{thm:frame-offset} reads this ratio as a Lorentz factor between the matter and the
photon frame.  The source supplies no derivation of either input from the HC trace formula: the
proof of Theorem~\ref{thm:frame-offset} says that \(\gstring\) is computed in the photon frame
from that formula, while the proof of the OP~1/T4 statement around
Remark~\ref{rem:goe-gue-string-lqg} calls it extrinsically given.  The two statements are
incompatible, and only the second is supported by the text.

\paragraph{F11a --- the entropy factor.}
The remark ``The entropy factor'' in \S\ref{sec:bh} writes the LQG entropy as
\(S=A\ln2/(8\pi\gBI\lP^2)\), finds agreement with \(S_{BH}=A/(4\lP^2)\) at
\(\gBI=\ln2/(2\pi)\approx0.110\), and records a discrepancy factor \(2.15\) as an open quantitative
problem.  The entropy formula of the counting that produced \(\ln2/(\pi\sqrt3)\) is
\[
S=\frac{\ln2}{4\pi\sqrt3\,\gamma}\,\frac{A}{\lP^2},
\]
which agrees with \(S_{BH}\) exactly at \(\gamma=\ln2/(\pi\sqrt3)=\gstring\).  The factor \(2.15\)
is produced by the formula used in the remark and is not a property of the framework; the
discrepancy between the two countings is \(\gBI/\gstring=\gL\) itself.

\paragraph{Consequence.}
The agreement of \(1-\bsym^2\) with the measured \(\Omega_\Lambda\) at the \(1\)--\(2\%\) level
recorded in \S\ref{subsec:v480-r03-p4-recheck} rests on \(\bsym\), and \(\bsym\) rests on the ratio
of two countings of one parameter, one of which is superseded.  As long as that is not stated,
the agreement reads as better supported than it is.

\paragraph{A reading in the readout vocabulary, offered and not adopted.}
The superseded counting admits only minimal spins \(j=\tfrac12\) at the horizon punctures, the
corrected one all spins.  \(\gL\) could therefore be read as the ratio between a minimal-spin
readout and a full-spin readout of the same horizon, two readouts in the sense of r81--r82 and
r93.  Under that reading \(\gL\) is not a Lorentz factor, and the chain
\(\gL\to\bsym\to\Omega_\Lambda\) would need a new justification: why the cosmological fraction
should be governed by the ratio of two horizon readouts.  r03 records the reading as an option
and asserts nothing beyond the provenance above.  Whether \(\gstring\) is renamed, the chain is
rederived, or this reading is adopted is left to the author.

\subsection*{v4.8.0-r03 status of the affected statements}
\label{subsec:v480-r03-status}

\[
\begin{array}{l|l|l}
\text{statement} & \text{finding} & \text{status after r03}\\
\hline
\text{Eq.~\eqref{eq:rho-ratio}, Cor.~\ref{cor:matter-error}} & \text{F1} & \text{corrected: } \gL^5\approx22.5\\
\alpha\text{-drift corollary} & \text{F2} & \text{exponent corrected; no observed drift}\\
\Lambda\text{ in Prop.~\ref{prop:bare-lambda}, Thm.~\ref{thm:dark-offset}, P4} & \text{F3} & \text{open: one type for }\Lambda\text{ to be chosen}\\
\text{Rem.~\ref{rem:open}, P4} & \text{F4} & \text{contradicted as stated; }\Omega\text{-reading offered}\\
\text{Thm.~\ref{thm:lambda-zero}(iii)} & \text{F8, F9} & \text{value off by }10^{121}\text{; citation corrected}\\
\text{Eq.~\eqref{eq:tr-balance}} & \text{F10} & \text{Planck units made explicit}\\
\gstring,\ \text{Thm.~\ref{thm:frame-offset}} & \text{F11} & \text{provenance recorded; readout reading offered}\\
\text{entropy-factor remark, \S\ref{sec:bh}} & \text{F11a} & \text{corrected: no factor }2.15\\
\Omega\text{ statements, Rem.~\ref{rem:omega-bsym}} & \text{F5} & \text{values corrected; two predictions conflict}\\
\text{Prop.~\ref{prop:hawking-gBI-tachyon}} & \text{F6} & \text{assumption } k(M)=0 \text{ made explicit}\\
\text{Thm.~\ref{thm:dark-offset}} & \text{F7} & \text{notation only}
\end{array}
\]

None of these statements is one of the guarded claims of the release scope, and none of the
five guards depends on them.  F3 and F4 concern the photon-frame \(\Lambda\) programme, which the
source already marks as open; r03 sharpens that status.  For P4 the source-only layer records
the contradiction and does not supply a replacement; a reformulation of the
\(\Lambda\) prediction, or its withdrawal, is left to the author.

\subsection*{v4.8.0-r03 ledger and guard update}
\label{subsec:v480-r03-ledger-guard}

\[
R_{\rm dim}^{v480\text{-}r03}=(0,0,1,1,1,1)
\]
for (\(k\)-calculus corrections F1, F2, F6 derived; numerical corrections F1, F5 recorded; single
type for \(\Lambda\) not chosen; a replacement for P4 offered in \S\ref{subsec:v480-r03-p4-recheck} but not adopted; one of the two
\(\Omega_\Lambda/\Omega_M\) predictions not selected; independent justification of \(\gL\) as a frame factor not supplied, \S\ref{subsec:v480-r03-frame-parameter-provenance}).

The layer spends no guard:
\[
R_{\rm glob}^{(v480\text{-}r03,+)}=(0,1,1,1,1),
\qquad
\Gamma_{\rm guard}^{v480\text{-}r03}=(1,1,1,1,1),
\qquad
\beta_4=1/256\ \text{unspent}.
\]
It does not imply exact finite \(C5\) zero closure, exact all-readout terminality,
\(C6/T2\), a global photon-clock/RH theorem, or global all-readout closure.

% ============================================================================
\clearpage
\section*{v4.8.0-r04 Source-Only: Decoupling from the Riemann Source and Standalone Sphaera Transformations}
\addcontentsline{toc}{section}{v4.8.0-r04 Source-Only: Decoupling from the Riemann Source and Standalone Sphaera Transformations}
\label{sec:v480-r04-riemann-decoupling}
% ============================================================================

\paragraph{Purpose.}
The Companion cites the Riemann source \cite{brendecke2026rh} at 23 sites for seven distinct
results.  A full audit of that source (v4.15) and a recomputation of its central formulas
show that its proof chain for RH does not hold.  The author used the source as a starting
point; what the Companion needs from it are the Heisenberg compensator and the Sphaera model
of the basic transformations (reflection, cascade, triole hierarchy).  This layer records the
audit, gives standalone definitions of those three objects that do not depend on any RH
statement, repairs the triole hierarchy where it fails, and reclassifies every other
dependency.  Earlier layers are not rewritten.

\noindent\textbf{Source-only status.}
This v4.8.0-r04 layer is source-only.  It does not build a PDF, does not assign a new DOI,
and does not authorize a public release.  The v4.8.0 Version DOI
\href{https://doi.org/10.5281/zenodo.23365244}{10.5281/zenodo.23365244} and the Concept DOI
\href{https://doi.org/10.5281/zenodo.19210728}{10.5281/zenodo.19210728} are unchanged by it.

\subsection*{v4.8.0-r04 audit of the Riemann source}
\label{subsec:v480-r04-source-audit}

The combinatorial orbit theory (L1--L11, Axioms W and G, K1), the arithmetic identification
\(\mathcal Z=\zeta\) for \(\operatorname{Re}s>1\) and the functional-equation sketch are correct.
The following findings concern the steps that carry the claim.  Numerical checks are recorded
in the audit scripts of the working folder \texttt{llm/}.

\paragraph{A1 --- explicit formula.}
The source states \(\Theta(\tau)=1-\Sigma(\tau)+R(\tau)\) with
\(\Theta=\sum_{p,m}m^{-1}\delta(\tau-m\log p)\).  Against a smooth test function supported in
\((0.6,3.2)\) with 300 zeros the two sides give \(2.1587\) and \(0.5198\).  The correct form,
which matches to six digits (\(4.393268\) on both sides), is
\[
\tau\,\Theta(\tau)=\sum_{n}\Lambda(n)\,\delta(\tau-\log n)
=e^{\tau}-\sum_{\rho}e^{\rho\tau}-\sum_{k\ge1}e^{-2k\tau}\qquad(\tau>0).
\]

\paragraph{A2 --- the reflection on the Hardy space.}
\(s\mapsto1-\bar s\) maps \(\operatorname{Re}s>\tfrac12\) outside itself.  It acts only on
boundary values, where it is complex conjugation; an \(H^2\) function with real boundary
values vanishes.  In the only well-defined reading the source's \(\Hplusplus\) is \(\{0\}\).
The antilinearity of \(J\) also makes \(\tfrac12(1\pm J)\) only real-linear.

\paragraph{A3 --- self-adjointness.}
A conjugation commuting with a symmetric operator gives equal deficiency indices, not
essential self-adjointness.  The operator \(-i(x\,d/dx+\tfrac12)\) is already self-adjoint on
\(L^2(\mathbb R^+,dx)\), with continuous spectrum \(\mathbb R\); self-adjointness is not the
bottleneck.

\paragraph{A4 --- discreteness and identification.}
The one-dimensional relation \([Q,P]=i\) does not produce a representation of an adelic Weyl
algebra; Stone--von Neumann for a locally compact abelian group yields \(L^2\) of the group,
not of \(\mathbb Q^+\backslash\mathbb A^1\).  The identification step uses Connes' trace formula
on the adele class space, whose validity is equivalent to RH; moreover the source's own formula
gives \(\mathcal L[\Sigma](s)=1/s-\log\zeta(s)+\mathcal L[R](s)\), not \(\log\zeta(s)\), and
\(\mathcal L[R]\) diverges.

\paragraph{A5 --- the set argument.}
The source assigns the same kernel \((\mathbb P,\mathrm{id},\log p)\) to \(\zeta\) and to every
self-dual automorphic \(L\)-function, and its operator stream is independent of the
\(L\)-function.  It would therefore assign one spectrum to all of them, while their zeros differ:
the first zero of \(\zeta\) is \(\tfrac12+14.1347251417\,i\), that of \(L(s,\chi_{-4})\) is
\(\tfrac12+6.0209489047\,i\).  The Davenport--Heilbronn function satisfies a Riemann-type
functional equation and has the zero \(0.8085171825+85.69934849\,i\) off the critical line.
Where the zeros lie is therefore not determined by the set \((\mathbb P,\log p)\) together with
the functional equation; it requires the multiplicative structure of the Euler product.

\paragraph{Theoremlet v4.8.0-r04.1.}
The Riemann source is not a proof source for RH.  None of the five guards of the release scope
depends on it, and the guard \(W_\gamma^{\rm glob}\), which forbids a global photon-clock/RH
theorem, is confirmed rather than affected.

\paragraph{Proof sketch.}
A1--A4 each break a link of the chain K1/L11 \(\to\) T1+T3 \(\to\) T4+T5 \(\to\) T6, and A5 shows
that the chain could not distinguish \(\zeta\) from other \(L\)-functions.  The guards concern
\(C5\), all-readout terminality, \(C6/T2\), \(W_\gamma^{\rm glob}\) and \(W_{\rm all}\); the
Companion already treated the RH channel as a local photon-clock probe, so no guarded statement
used RH as an input.  \(\square\)

\subsection*{v4.8.0-r04 standalone Heisenberg compensator}
\label{subsec:v480-r04-standalone-compensator}

\paragraph{Definition v4.8.0-r04.1.}
Let \(\mathcal H\) be a complex Hilbert space and \(J:\mathcal H\to\mathcal H\) a
\emph{linear} unitary involution, \(J^2=1\), \(J^*=J\).  The Heisenberg compensator and its
antisymmetric partner are
\[
C:=\tfrac12(1+J),\qquad A:=\tfrac12(1-J).
\]

\paragraph{Theoremlet v4.8.0-r04.2.}
\(C\) and \(A\) are complementary orthogonal projections, \(C+A=1\), \(CA=0\);
\(\operatorname{ran}C=\ker(J-1)\), \(\operatorname{ran}A=\ker(J+1)\); \(C(Jf)=C(f)\) and
\(A(Jf)=-A(f)\).  The terms \(\mathcal L_{\rm HC}\) and \(\mathcal L_{\rm frame}\) of
Theorem~\ref{thm:gauge-invariance} are \(J\)-invariant.  \(\mathcal L_{\rm kin}\) is
\(J\)-invariant if and only if \(JL_0J=L_0\).

\paragraph{Proof sketch.}
\(C^2=\tfrac14(1+2J+J^2)=C\) and \(C^*=C\) since \(J^*=J\); likewise for \(A\), and
\(CA=\tfrac14(1-J^2)=0\).  \(CJ=\tfrac12(J+1)=C\) and \(AJ=-A\).  \(\langle Jf,Jf\rangle=\langle
f,f\rangle\) by unitarity; \(\langle Jf,L_0Jf\rangle=\langle f,JL_0Jf\rangle\).  \(\square\)

\paragraph{Instances.}
(a) The \(e_1\)-reflection on \(\mathbb C\otimes\HH\) with the Euclidean inner product of the
coefficients is a linear unitary involution; on \(\HH\) it is an anti-automorphism.
(b) The linear reflection \(s\mapsto1-s\) restricted to the critical line acts on the ordinate as
the parity \(t\mapsto-t\), \((Jf)(t)=f(-t)\) on \(L^2(\mathbb R)\).  For (b), in logarithmic
coordinates \(L_0=-i(\partial_y+\tfrac12)\) and parity sends \(\partial_y\mapsto-\partial_y\), so
\(JL_0J\neq L_0\): the kinetic invariance is open for this instance.

The compensator therefore requires no RH statement.  What it loses relative to the earlier
text is the claim that its image is a space whose \(L_0\)-spectrum consists of zeta zeros.

\subsection*{v4.8.0-r04 standalone Sphaera cascade}
\label{subsec:v480-r04-standalone-cascade}

\paragraph{Definition v4.8.0-r04.2.}
Let \(\mathbb S\) be a fixed Hilbert space (a set, not a universal collection: ZF admits no
set of all sets), and let \(\mathbb S=S_N\supseteq S_{N-1}\supseteq\cdots\supseteq S_0\) be
closed subspaces with orthogonal projections \(\Pi_n:S_n\to S_{n-1}\).  The cascade is
\(\KC:=\Pi_1\circ\cdots\circ\Pi_N\).

\paragraph{Theoremlet v4.8.0-r04.3.}
\(\KC\) is the orthogonal projection of \(\mathbb S\) onto \(S_0\); in particular
\(\KC^2=\KC\), and \(\KC f=f\) if and only if \(f\in S_0\).

\paragraph{Proof sketch.}
For nested closed subspaces \(V\subseteq W\) the projections satisfy \(P_VP_W=P_V\); induction
over the chain gives \(\KC=P_{S_0}\).  \(\square\)

The Maximal Totality principle is read inside this definition: \(\KC\) is defined on
\(\mathbb S\) and maps every element of \(\mathbb S\) into \(S_0\).  It is not applied to
elements outside its domain.

\subsection*{v4.8.0-r04 repaired triole hierarchy}
\label{subsec:v480-r04-triole-hierarchy}

Use the Cayley--Dickson convention of the Riemann source,
\((p+q\ell)(r+s\ell)=(pr-\bar sq)+(sp+q\bar r)\ell\), with \(e_4=\ell\), \(e_5=e_1\ell\),
\(e_6=e_2\ell\), \(e_7=e_3\ell\).  Explicit multiplication gives the seven quaternionic triples
(lines of the Fano plane)
\[
(1,2,3),\ (1,4,5),\ (1,6,7),\ (2,4,6),\ (2,5,7),\ (3,4,7),\ (3,5,6).
\]

\paragraph{Theoremlet v4.8.0-r04.4.}
Three quaternionic triples cover all seven imaginary units of \(\OO\) if and only if they pass
through a common unit.  There are exactly seven such configurations, one per unit.  The source's
second triole \((e_4,e_5,e_6)\) is not quaternionic, since \(e_4e_5=e_1\).

\paragraph{Proof sketch.}
Two distinct lines of the Fano plane meet in exactly one point.  For three distinct lines,
\(|L_1\cup L_2\cup L_3|=9-|L_1\cap L_2|-|L_1\cap L_3|-|L_2\cap L_3|+|L_1\cap L_2\cap L_3|\),
which is \(7\) if the three lines are concurrent and \(6\) otherwise.  Each point lies on three
lines, so each point gives exactly one concurrent triple.  \(e_4e_5=\ell(e_1\ell)=e_1\) is the
source's own computation.  \(\square\)

\paragraph{Replacement.}
The triole configuration through \(e_4=\ell\) is
\[
\mathcal T_1'=(e_1,e_4,e_5),\qquad \mathcal T_2'=(e_2,e_4,e_6),\qquad
\mathcal T_3'=(e_3,e_4,e_7),
\]
which keeps the source's third triole and shares one common axis.  Equally admissible is the
configuration through \(e_1\), which keeps the inner triole \((e_1,e_2,e_3)\).

\paragraph{Reflection as algebra symmetry.}
The reflection negating \(e_1\) alone is not an automorphism of \(\OO\).  Of the
\(2^7\) sign changes exactly eight are automorphisms; those negating \(e_1\) negate four units
and fix one quaternionic line pointwise, for example \(\{e_1,e_2,e_4,e_7\}\) with fixed line
\((3,5,6)\).  A reflection meant as a symmetry of the algebra must have this form.

\paragraph{Above the octonions.}
In the sedenions the source's outer octo-triole \((f_8,\dots,f_{14})\) is not closed
(\(f_8f_9=f_1\)), and zero divisors exist, e.g. \((e_3+f_{10})(e_6-f_{15})=0\).  The hierarchy
continues algebraically, not geometrically; the Hopf structure ends at \(\OO\), as the source
itself states.  Theorem~\ref{thm:three-generations} uses \(e_4,e_5,e_6\) as an \(SU(3)\)
triplet, not as a quaternionic triole, and is not changed by this layer.

\subsection*{v4.8.0-r04 dependency map}
\label{subsec:v480-r04-dependency-map}

\begin{center}
\small
\begin{tabular}{@{}>{\raggedright\arraybackslash}p{0.31\textwidth}>{\raggedright\arraybackslash}p{0.30\textwidth}>{\raggedright\arraybackslash}p{0.33\textwidth}@{}}
\toprule
site & used from the Riemann source & status after r04\\
\midrule
Thm.~\ref{thm:unity}, equivalence clause & Stone--von Neumann (Lem.~13.2, Step~3) & open\\
Thm.~\ref{thm:unitary-equiv} & Stone--von Neumann, T4+T5 & open\\
Thms.~\ref{thm:bell-consistency}, \ref{thm:string-lqg-duality} & Stone--von Neumann & equivalence open\\
Thms.~\ref{thm:einstein}, \ref{thm:gue} & T6, \(\Sigma=1-\Theta+R\) & open; formula to be replaced (A1)\\
Thm.~\ref{thm:toe}, row ``Riemann Hypothesis'' & RH as a result & withdrawn: violates \(W_\gamma^{\rm glob}\)\\
Thm.~\ref{thm:gauge-invariance} & \(\mathrm{Adj}_J(L_0)=0\) & \(\mathcal L_{\rm HC}\), \(\mathcal L_{\rm frame}\) standalone; \(\mathcal L_{\rm kin}\) open\\
Thm.~\ref{thm:eom} & field equation & projector algebra kept; zero reading withdrawn\\
Thms.~\ref{thm:op2-existence}, \ref{thm:op2-selfadj} & T1+T3, T4+T5 & open (OP2 stays open)\\
Thm.~\ref{thm:partial-cancel}, Rem.~\ref{rem:explicit-formula-remainder} & explicit formula (Thm.~3.3) & to be replaced by A1\\
prime lattice \(\Gamma_P\) & classical arithmetic & unaffected\\
\S\ref{sec:rh} & RH as consequence & read as conjecture only\\
\bottomrule
\end{tabular}
\end{center}

\paragraph{Theoremlet v4.8.0-r04.5.}
After this layer the Heisenberg compensator, the Sphaera cascade and the triole hierarchy up to
\(\OO\) rest on Definitions v4.8.0-r04.1 and r04.2 and on Theoremlet r04.4, without any statement
of \cite{brendecke2026rh} beyond classical arithmetic.  The statements listed as open lose their
previous justification but are not refuted by this layer.

\paragraph{Proof sketch.}
Theoremlets r04.2--r04.4 use only Hilbert-space projections and explicit Cayley--Dickson
multiplication.  For the open items the audit A1--A5 removes the cited justification; whether
an independent justification exists is not decided here.  \(\square\)

\subsection*{v4.8.0-r04 ledger and guard update}
\label{subsec:v480-r04-ledger-guard}

\[
R_{\rm dec}^{v480\text{-}r04}=(0,0,0,1,1,1)
\]
for (standalone compensator; standalone cascade; repaired triole hierarchy; Stone--von Neumann
equivalence of the kernels open; trace-formula field equations open pending A1; kinetic
\(J\)-invariance open).

The layer spends no guard:
\[
R_{\rm glob}^{(v480\text{-}r04,+)}=(0,1,1,1,1),
\qquad
\Gamma_{\rm guard}^{v480\text{-}r04}=(1,1,1,1,1),
\qquad
\beta_4=1/256\ \text{unspent}.
\]
It does not imply exact finite \(C5\) zero closure, exact all-readout terminality,
\(C6/T2\), a global photon-clock/RH theorem, or global all-readout closure.


% ============================================================================
\clearpage
\section*{v4.8.0-r05 Source-Only: Corrected Explicit Formula in the Field-Equation and \texorpdfstring{$\Theta$}{Theta}--\texorpdfstring{$R$}{R} Layer}
\addcontentsline{toc}{section}{v4.8.0-r05 Source-Only: Corrected Explicit Formula in the Field-Equation and Theta--R Layer}
\label{sec:v480-r05-corrected-explicit-formula}
% ============================================================================

\paragraph{Purpose.}
Layer v4.8.0-r04 recorded that the explicit formula \(\Sigma=1-\Theta+R\) taken from
\cite{brendecke2026rh} is misstated (finding A1) and listed the sites that use it.  This layer
carries the corrected formula into those sites: the Einstein-equation reading of
Theorem~\ref{thm:einstein}, the bare-\(\Lambda\) proposition, the \(\Theta\)--\(R\) cancellation
statement, the photon-frame \(\Theta\)-sum conjecture of OP~5 Step~(B), and the GUE theorem.
Two statements can now be decided numerically; both are decided here.  Earlier layers are not
rewritten.

\noindent\textbf{Source-only status.}
This v4.8.0-r05 layer is source-only.  It does not build a PDF, does not assign a new DOI,
and does not authorize a public release.  The v4.8.0 Version DOI
\href{https://doi.org/10.5281/zenodo.23365244}{10.5281/zenodo.23365244} and the Concept DOI
\href{https://doi.org/10.5281/zenodo.19210728}{10.5281/zenodo.19210728} are unchanged by it.

\subsection*{v4.8.0-r05 corrected trace decomposition}
\label{subsec:v480-r05-trace-decomposition}

Write \(\Sigma(\tau):=\sum_\rho e^{i\gamma_\rho\tau}\), summed over all non-trivial zeros
\(\rho=\tfrac12+i\gamma_\rho\) (the zeros are used as numerically known data; no statement about
their real parts is assumed).

\paragraph{Theoremlet v4.8.0-r05.1.}
For \(\tau>0\), in the sense of distributions,
\[
\Sigma(\tau)
=\underbrace{e^{\tau/2}}_{M(\tau)}
\;-\;\underbrace{e^{-\tau/2}\sum_{n\ge2}\Lambda(n)\,\delta(\tau-\log n)}_{P(\tau)}
\;-\;\underbrace{\frac{e^{-\tau/2}}{e^{2\tau}-1}}_{Z(\tau)} .
\]
\(M\) comes from the pole of \(\zeta\) at \(s=1\), \(P\) from the prime powers with von Mangoldt
weights, \(Z\) from the trivial zeros.

\paragraph{Proof sketch.}
This is the classical explicit formula
\(\sum_n\Lambda(n)\delta(\tau-\log n)=e^\tau-\sum_\rho e^{\rho\tau}-\sum_{k\ge1}e^{-2k\tau}\),
multiplied by \(e^{-\tau/2}\), with \(\sum_k e^{-2k\tau}=1/(e^{2\tau}-1)\).  Paired with a smooth
test function supported in \((0.6,3.2)\) and 300 zeros, both sides of the unmultiplied identity
agree to six digits (\(4.393268\)), while the form \(1-\Theta+R\) of the source gives
\(2.1587\) against \(0.5198\) (r04, A1).  \(\square\)

\paragraph{What changes.}
The three-part structure survives: pole, primes, trivial zeros.  What changes is every
functional form.  The constant \(1\) becomes the growing term \(e^{\tau/2}\); the primes enter with
weight \(\Lambda(n)=\log p\) and damping \(e^{-\tau/2}\), not with weight \(1/m\); the trivial-zero
term has denominator \(e^{2\tau}-1\) and positive sign inside \(Z\).  In addition,
equation~\eqref{eq:Theta-def} uses \(\Theta\) for the zero sum itself, which is also the
definition of \(\Sigma\); read literally, \eqref{eq:gw} then gives \(2\Sigma=1+R\).  The symbol
\(\Theta\) is reserved here for the prime-side distribution only.

\subsection*{v4.8.0-r05 the Einstein-equation reading}
\label{subsec:v480-r05-einstein-reading}

Theorem~\ref{thm:einstein} varies \(S=\int_0^\infty\Sigma(\tau)\,d\tau\).  With the corrected
decomposition this integral does not exist: \(M\) grows like \(e^{\tau/2}\), \(Z(\tau)\sim1/(2\tau)\)
as \(\tau\to0^+\), and \(\Sigma\) is a distribution, not a function.  Moreover
\(\mathrm{Tr}\,e^{-i\tau L_0}\) is the trace of a unitary group, not a heat kernel, so the heat-kernel
variation formula used in the proof does not apply.

\paragraph{Theoremlet v4.8.0-r05.2.}
The correspondence
\[
M\ \leftrightarrow\ \Lambda g_{\mu\nu},\qquad
P\ \leftrightarrow\ T_{\mu\nu},\qquad
Z\ \leftrightarrow\ G_{\mu\nu}
\]
is retained as a structural conjecture.  Theorem~\ref{thm:einstein} is not derived: a well-defined
action requires a test-function pairing \(\langle\Sigma,\varphi\rangle\) and an explicit dependence
of \(\varphi\) on \(g^{\mu\nu}\), neither of which is supplied.

\paragraph{Proof sketch.}
Divergence of the three pieces is read off from Theoremlet r05.1; the trace statement follows
from the unitarity of \(e^{-i\tau L_0}\).  \(\square\)

\paragraph{Bare \(\Lambda\).}
Proposition~\ref{prop:bare-lambda} assigns to the constant term a size \(\gamma_1^2\).  In the corrected
formula the pole term is \(M(\tau)=e^{\tau/2}\); it carries no dependence on any zero, so no
\(\gamma_1^2\) can arise from it.  The proposition is withdrawn as a derivation; the scale
\(\gamma_1^2\) remains an unsupported ansatz.

\paragraph{Cancellation.}
Theorem~\ref{thm:partial-cancel} reads the explicit formula as a cancellation of the pole by the
other terms.  The formula is an identity between distributions; it does not by itself fix the value
of a truncated integral.  The open content of Step~(B) is exactly such a value, decided below.

\subsection*{v4.8.0-r05 decision of the \texorpdfstring{$\Theta$}{Theta}-sum conjecture}
\label{subsec:v480-r05-theta-sum-decision}

Proposition~\ref{conj:theta-sum-B} states
\[
\mathcal T(T)=2\sum_{n\ge1}\frac{\sin(T\gamma_n/\gL)}{\gamma_n/\gL}
=\frac{3T}{2}-\log T-\frac{T^2}{24}-\gamma_E-\bsym^2T^2\approx0.8934,
\qquad T=\gL^{-3/2}\approx0.3928 .
\]

\paragraph{Theoremlet v4.8.0-r05.3.}
\(\mathcal T(T)=-0.1097\pm0.0009\).  The conjecture is false; the two sides differ by \(1.003\).
Satisfying the equation would require \(\bsym^2\approx6.80\), outside \([0,1)\).

\paragraph{Proof sketch.}
\(\mathcal T(T)=\gL\,S(a)\) with \(S(a)=2\sum_{\gamma>0}\sin(a\gamma)/\gamma\) and
\(a=T/\gL\approx0.2107<\log2\).  For \(x=e^a\in(1,2)\) one has \(\psi_0(x)=0\), so the
Riemann--von Mangoldt formula gives exactly
\(\sum_\rho x^\rho/\rho=x-\log2\pi-\tfrac12\log(1-x^{-2})\).  Pairing \(\rho\) with \(\bar\rho\),
\[
S(a)=e^{-a/2}\Bigl(x-\log2\pi-\tfrac12\log(1-x^{-2})\Bigr)
-\sum_{\gamma>0}\frac{\cos a\gamma}{\tfrac14+\gamma^2}
+\sum_{\gamma>0}\frac{\sin a\gamma}{2\gamma(\tfrac14+\gamma^2)},
\]
with both remaining sums absolutely convergent.  With the first 2000 zeros the value is
\(-0.109652\) (50 zeros: \(-0.109464\); 500: \(-0.109661\)), and the tail is bounded by
\(8.2\cdot10^{-4}\) through the zero density \((2\pi)^{-1}\log(\gamma/2\pi)\).  The direct partial sums
converge to the same value (\(-0.1449\), \(-0.1117\), \(-0.1053\) at \(N=50,500,2000\)); the
first agrees with the \(-0.146\) recorded in Remark~\ref{rem:theta-sum-structure}.  \(\square\)

\paragraph{Consequence.}
By Remark~\ref{rem:explicit-formula-remainder} the conjecture was stated as equivalent to the
balance condition~\eqref{eq:tr-balance} and to OP~5 Step~(B).  That equivalence used the old
\(R\); with the corrected decomposition it would have to be re-derived.  Either way, the balance
statement in the form given there does not hold.  The defect of almost exactly \(1\) suggests that
the constant term was mishandled in the regularization; even a shift by \(1\) leaves a gap of
\(0.003\), larger than the error bound.  Prediction P4 and the photon-frame \(\Lambda\) programme,
already marked open in r03, lose this route.

\subsection*{v4.8.0-r05 GUE statistics: empirical, not derived}
\label{subsec:v480-r05-gue}

Theorem~\ref{thm:gue} derives GUE statistics from self-adjointness, irreducibility and the
absence of time reversal.  The derivation does not hold: self-adjointness and irreducibility do not
imply any random-matrix class (the harmonic oscillator has neither GUE nor GOE spacings), and an
antiunitary \(J\) with \(J^2=1\) commuting with \(L_0\), as asserted in the proof, would be a
time-reversal-type symmetry and would point to GOE, not GUE.

\paragraph{Theoremlet v4.8.0-r05.4.}
The nearest-neighbour spacings of the zeros follow GUE statistics empirically.  For zeros
\(101\)--\(2000\), unfolded with \(N(t)=\tfrac{t}{2\pi}\log\tfrac{t}{2\pi e}+\tfrac78\), the
Kolmogorov--Smirnov distance to the Wigner surmise is \(0.041\) for GUE (\(\beta=2\)), \(0.108\)
for GOE (\(\beta=1\)) and \(0.315\) for Poisson; small spacings \(s<0.1\) occur with frequency
\(0.0005\) (GUE surmise \(0.0011\), GOE \(0.0078\), Poisson \(0.095\)).

\paragraph{Proof sketch.}
Direct computation from 2000 zeros.  The residual GUE distance reflects the Wigner surmise as an
approximation and the low height; it is consistent with the Montgomery--Odlyzko results.
\(\square\)

Theorem~\ref{thm:gue} is therefore retained as an empirical statement about the zeros and
withdrawn as a consequence of the Sphaera structure.

\subsection*{v4.8.0-r05 status of the affected statements}
\label{subsec:v480-r05-status}

\begin{center}
\small
\begin{tabular}{@{}>{\raggedright\arraybackslash}p{0.34\textwidth}>{\raggedright\arraybackslash}p{0.58\textwidth}@{}}
\toprule
statement & status after r05\\
\midrule
\eqref{eq:gw}, \eqref{eq:Theta-def}, \eqref{eq:R-def} & replaced by Theoremlet r05.1\\
Thm.~\ref{thm:einstein} & structural conjecture (r05.2); action not defined\\
Prop.~\ref{prop:bare-lambda} & withdrawn as derivation; \(\gamma_1^2\) unsupported\\
Thm.~\ref{thm:partial-cancel} & identity, not cancellation; content moved to Step~(B)\\
Prop.~\ref{conj:theta-sum-B}, Rem.~\ref{rem:explicit-formula-remainder} & false: \(\mathcal T(T)=-0.1097\), not \(0.893\)\\
Thm.~\ref{thm:gue} & empirical fact (r05.4); derivation withdrawn\\
\bottomrule
\end{tabular}
\end{center}

\subsection*{v4.8.0-r05 ledger and guard update}
\label{subsec:v480-r05-ledger-guard}

\[
R_{\rm fe}^{v480\text{-}r05}=(0,0,1,1,0)
\]
for (corrected trace decomposition; \(\Theta\)-sum conjecture decided; Einstein-equation reading
not derived; photon-frame balance condition without a valid form; GUE statement settled
empirically).

The layer spends no guard:
\[
R_{\rm glob}^{(v480\text{-}r05,+)}=(0,1,1,1,1),
\qquad
\Gamma_{\rm guard}^{v480\text{-}r05}=(1,1,1,1,1),
\qquad
\beta_4=1/256\ \text{unspent}.
\]
It does not imply exact finite \(C5\) zero closure, exact all-readout terminality,
\(C6/T2\), a global photon-clock/RH theorem, or global all-readout closure.


% ============================================================================
\clearpage
\section*{v4.8.0-r06 Source-Only: Signature of the Biquaternion Kernel and the Complex-Vector Reading}
\addcontentsline{toc}{section}{v4.8.0-r06 Source-Only: Signature of the Biquaternion Kernel and the Complex-Vector Reading}
\label{sec:v480-r06-complex-vector-reading}
% ============================================================================

\paragraph{Purpose.}
Definition~\ref{def:biquaternion} gives the kernel
\(q_n=it\cdot\mathbf 1+z_1e_1+z_2e_2+z_3e_3\) with \(z_k\in\mathbb C\) and the norm
\(|q_n|^2=-(ct)^2+|z_1|^2+|z_2|^2+|z_3|^2\), read as a Minkowski metric of signature
\((-,+,+,+)\) on which the Lorentz group acts by quaternion conjugation.  This layer computes
the signature, identifies the natural invariant of \(\mathbb C\otimes\HH\), and adopts as working
reading the one under which the kernels \((z_1,z_2,z_3)\) carry the Lorentz group.  In that reading
the symmetric triole is exactly the light-like centred configuration, and the photon sector is the
base \(S_0\cong\mathbb CP^1\) of Theorem~\ref{thm:hopf}.  Earlier layers are not rewritten.

\noindent\textbf{Source-only status.}
This v4.8.0-r06 layer is source-only.  It does not build a PDF, does not assign a new DOI,
and does not authorize a public release.  The v4.8.0 Version DOI
\href{https://doi.org/10.5281/zenodo.23365244}{10.5281/zenodo.23365244} and the Concept DOI
\href{https://doi.org/10.5281/zenodo.19210728}{10.5281/zenodo.19210728} are unchanged by it.

\subsection*{v4.8.0-r06 signature and natural norm}
\label{subsec:v480-r06-signature}

Write \(z_k=x_k+iy_k\) with \(x_k,y_k\in\mathbb R\).

\paragraph{Theoremlet v4.8.0-r06.1.}
\begin{enumerate}[label=(\roman*)]
\item The norm of Definition~\ref{def:biquaternion} is the real quadratic form
\(-t^2+\sum_k(x_k^2+y_k^2)\) of signature \((1,6)\), not \((1,3)\).
\item Conjugation \(q\mapsto uqu^{-1}\) by \(u\in\HH^\times\) fixes the scalar part and acts by
\(SO(3)\) on \((x_k)\) and on \((y_k)\) separately; it contains no boost.
\item The natural invariant of \(\mathbb C\otimes\HH\) is the complex-bilinear norm
\[
N(q):=q\bar q=-t^2+z_1^2+z_2^2+z_3^2,\qquad
\operatorname{Re}N=-t^2+|x|^2-|y|^2,\quad \operatorname{Im}N=2\,x\cdot y,
\]
where \(\bar q\) is quaternion conjugation without complex conjugation; \(N(pq)=N(p)N(q)\).
\end{enumerate}

\paragraph{Proof sketch.}
(i) The Hessian of the form has eigenvalues \(-1\) once and \(+1\) six times.  (ii) For real
\(u\), \(u e_k u^{-1}\) is a rotation of the imaginary units and \(u\mathbf 1u^{-1}=\mathbf 1\).
(iii) Direct expansion of \(q\bar q\); the vector part vanishes because the coefficients commute.
Multiplicativity is inherited from \(\HH\).  \(\square\)

Hence the Minkowski reading of Definition~\ref{def:biquaternion} and the Lorentz statement of the
subsequent remark do not hold as written, and Theorem~\ref{thm:lorentz} establishes invariance
under the involution \(J\), not under the Lorentz group.

\subsection*{v4.8.0-r06 two admissible readings}
\label{subsec:v480-r06-readings}

\paragraph{Reading A (four-vector).}
On the real slice \(\{q:q^\dagger=-q\}\), where \(q^\dagger\) combines quaternion and complex
conjugation, the \(z_k\) are real and \(N(q)=-t^2+x_1^2+x_2^2+x_3^2\) is exactly Minkowski
\((1,3)\).  The map \(q\mapsto uqu^\dagger\) with \(N(u)=1\) preserves the slice and \(N\), and
changes the time component: these are genuine boosts.  The kernels then carry no phase, so the
\(120^\circ\) triole has no place.

\paragraph{Reading B (complex three-vector).}
The kernels form \(z=(z_1,z_2,z_3)\in\mathbb C^3\), and \(q\mapsto uqu^{-1}\) with
\(u\in\mathbb C\otimes\HH\), \(N(u)=1\), fixes the scalar part and acts on \(z\) by
\(SO(3,\mathbb C)\).  Since \(SO(3,\mathbb C)\cong PSL(2,\mathbb C)\cong SO^+(1,3)\), this is the
proper orthochronous Lorentz group.  The invariant is \(s(z):=z_1^2+z_2^2+z_3^2\).  The conjugation
formula of the remark after Definition~\ref{def:biquaternion} becomes correct once \(u\) is taken
complex.  Formally, \(z\) is the Riemann--Silberstein vector \(E+iB\), and
\(s=E^2-B^2+2i\,E\cdot B\) collects the two field invariants.

\paragraph{Decision.}
Reading~B is adopted as the working reading of the kernel (author's decision, 2026-10-19).
Reading~A remains available as the four-vector structure of \(\mathbb C\otimes\HH\); the two
are not exclusive.  In Reading~B the scalar \(t\) is an invariant parameter, not the time
coordinate of a four-vector; its role is open.

\subsection*{v4.8.0-r06 the triole is the light-like centred configuration}
\label{subsec:v480-r06-triole-lemma}

\paragraph{Theoremlet v4.8.0-r06.2.}
For \((z_1,z_2,z_3)\in\mathbb C^3\),
\[
z_1+z_2+z_3=0\ \text{ and }\ z_1^2+z_2^2+z_3^2=0
\quad\Longleftrightarrow\quad
(z_1,z_2,z_3)=c\,(1,\omega,\omega^2)\ \text{up to order},
\]
with \(\omega=e^{2\pi i/3}\) and \(c\in\mathbb C\).

\paragraph{Proof sketch.}
With elementary symmetric functions \(e_1,e_2,e_3\) and \(p_2=\sum z_k^2\) one has
\(p_2=e_1^2-2e_2\).  Hence \(e_1=p_2=0\) gives \(e_2=0\), and the \(z_k\) are the roots of
\(\lambda^3-e_3\), i.e.\ the three cube roots of \(e_3\).  Conversely \(1+\omega+\omega^2=0\) and
\(1+\omega^2+\omega^4=0\).  \(\square\)

The triole condition of Theorem~\ref{thm:unity}, read with equal moduli and \(120^\circ\) phases,
is therefore not only a phase condition: in Reading~B it is equivalent to the Lorentz-invariant
null condition \(s=0\) together with centring \(\sum z_k=0\).

\subsection*{v4.8.0-r06 invariant sectors}
\label{subsec:v480-r06-sectors}

\paragraph{Theoremlet v4.8.0-r06.3.}
The orbits of \(SO(3,\mathbb C)\) on \(\mathbb C^3\) are the level sets \(\{s=c\}\) for \(c\neq0\),
the punctured null cone \(\{s=0\}\setminus\{0\}\), and \(\{0\}\).  The duality rotation
\(z\mapsto iz\) maps \(s\mapsto-s\) and preserves the null cone.

\paragraph{Proof sketch.}
For \(s\neq0\) put \(e_1=z/\sqrt s\), complete it to a complex-orthonormal basis with determinant
\(1\), and rotate \(z\) to \((\sqrt s,0,0)\); this was checked constructively.  Null vectors are
treated through Theoremlet r06.4 below.  \(s(iz)=-s(z)\) is immediate.  \(\square\)

\paragraph{Sector reading (conjecture).}
The classification of \S\ref{sec:unitary} by the sign of \(\operatorname{Re}z\) is not
Lorentz-invariant in Reading~B.  A Lorentz-invariant candidate is: \(s=0\) null (photon-type);
\(s\) real positive, electric-dominated; \(s\) real negative, magnetic-dominated.  The duality
\(z\mapsto iz\) exchanges the two massive types and fixes the null sector, matching the reading of
the tachyon as \(J\)-dual of matter and of the photon as fix-sector.  The assignment
electric \(\leftrightarrow\) matter, magnetic \(\leftrightarrow\) tachyon is a model decision and is
not made here; the meaning of non-real \(s\) is open.

\subsection*{v4.8.0-r06 the null cone and the photon base}
\label{subsec:v480-r06-null-cone}

For \(\psi=(\psi_1,\psi_2)\in\mathbb C^2\) put
\(z(\psi):=(\psi_1^2-\psi_2^2,\ i(\psi_1^2+\psi_2^2),\ -2\psi_1\psi_2)\).

\paragraph{Theoremlet v4.8.0-r06.4.}
\(z(\psi)\) lies on the null cone, every non-zero null vector has exactly the two preimages
\(\pm\psi\), and the projectivized null cone \(\{[z]\in\mathbb CP^2:s(z)=0\}\) is a smooth conic
isomorphic to \(\mathbb CP^1\cong S^2\).  Under the identification of Theorem~\ref{thm:hopf} this is
the photon base \(S_0\).  The map is equivariant: \(SL(2,\mathbb C)\) on \(\psi\) induces
\(SO(3,\mathbb C)\) on \(z\).

\paragraph{Proof sketch.}
\(s(z(\psi))=(a-b)^2-(a+b)^2+4ab=0\) with \(a=\psi_1^2\), \(b=\psi_2^2\).  For a null \(z\), set
\(\psi_1^2=\tfrac12(z_1-iz_2)\), \(\psi_2^2=\tfrac12(-z_1-iz_2)\); the null condition gives
\(z_3^2=4\psi_1^2\psi_2^2\), which fixes the relative sign, so the preimage is \(\pm\psi\).  The
conic is the Veronese image of \(\mathbb CP^1\).  Equivariance and invariance of the null cone
were checked numerically for random unit biquaternions.  \(\square\)

\paragraph{Consequences.}
The symmetric triole \((1,\omega,\omega^2)\) is the image of the spinor
\(\psi\approx(0.9744+0.1283i,\ 0.3097+0.4036i)\).  For the null axis of
Theorem~\ref{thm:fractal}, the natural Lorentz-invariant definition in Reading~B is \(s=0\), not
\(|q_n|^2=0\); its directions then form exactly \(S_0\).  The transport statement
\(\KC(L_n)=L_{n-1}\) of that theorem is not addressed and remains as stated there.

\paragraph{Information-level reading (conjecture).}
A normalized \(\psi\) is a qubit; its Bloch vector is the Hopf image \(S^3\to S^2\), and
\(k(\psi)=(|\psi|^2,\psi^\dagger\sigma\psi)\) is a light-like four-vector.  Every future time-like
vector is a sum \(k(\psi_a)+k(\psi_b)\) (checked constructively).  This supports the author's
conjecture that space-time is a projection of an information level, with single spinors as
light-like and spinor pairs as massive readouts, in line with
\(S_0=\mathbb CP^1\) and \(S_1=\HH P^1\) of Theorem~\ref{thm:hopf}.  No derivation of space-time
points is claimed.

\subsection*{v4.8.0-r06 status of the affected statements}
\label{subsec:v480-r06-status}

\begin{center}
\small
\begin{tabular}{@{}>{\raggedright\arraybackslash}p{0.34\textwidth}>{\raggedright\arraybackslash}p{0.58\textwidth}@{}}
\toprule
statement & status after r06\\
\midrule
Def.~\ref{def:biquaternion}, Minkowski reading & signature \((1,6)\); replaced by \(N(q)\) and Reading~B\\
Lorentz action by quaternion conjugation & correct only with complex \(u\), \(N(u)=1\)\\
Thm.~\ref{thm:lorentz} & shows \(J\)-invariance; Lorentz invariance supplied by r06 (Reading~B)\\
Sector classification, \S\ref{sec:unitary} & not invariant; invariant candidate via \(s\) (conjecture)\\
Triole condition, Thm.~\ref{thm:unity} & equivalent to centred null configuration (r06.2)\\
Null axis, Thm.~\ref{thm:fractal} & invariant definition \(s=0\); directions \(=S_0\) (r06.4)\\
Prop.~\ref{prop:orthogonality-third-kernel} & unaffected\\
\bottomrule
\end{tabular}
\end{center}

\subsection*{v4.8.0-r06 ledger and guard update}
\label{subsec:v480-r06-ledger-guard}

\[
R_{\rm kern}^{v480\text{-}r06}=(0,0,0,0,1,1)
\]
for (signature computed; Lorentz action supplied in Reading~B; triole lemma; null cone \(=S_0\);
sector assignment not chosen; role of \(t\) open).

The layer spends no guard:
\[
R_{\rm glob}^{(v480\text{-}r06,+)}=(0,1,1,1,1),
\qquad
\Gamma_{\rm guard}^{v480\text{-}r06}=(1,1,1,1,1),
\qquad
\beta_4=1/256\ \text{unspent}.
\]
It does not imply exact finite \(C5\) zero closure, exact all-readout terminality,
\(C6/T2\), a global photon-clock/RH theorem, or global all-readout closure.


% ============================================================================
\clearpage
\section*{v4.8.0-r07 Source-Only: Sector Assignment and Duality Angle in the Complex-Vector Reading}
\addcontentsline{toc}{section}{v4.8.0-r07 Source-Only: Sector Assignment and Duality Angle in the Complex-Vector Reading}
\label{sec:v480-r07-sector-assignment}
% ============================================================================

\paragraph{Purpose.}
Layer v4.8.0-r06 left the assignment of the two massive types of Reading~B to matter and tachyon
open.  This layer adopts a working assignment, justifies it from the signed tachyon-frame scale of
\S\ref{sec:r82-tachyon-frame-signed-scale}, and shows that every non-light-like kernel configuration
has a unique duality angle, so that non-real values of \(s\) need no separate sector.  Earlier
layers are not rewritten.

\noindent\textbf{Source-only status.}
This v4.8.0-r07 layer is source-only.  It does not build a PDF, does not assign a new DOI,
and does not authorize a public release.  The v4.8.0 Version DOI
\href{https://doi.org/10.5281/zenodo.23365244}{10.5281/zenodo.23365244} and the Concept DOI
\href{https://doi.org/10.5281/zenodo.19210728}{10.5281/zenodo.19210728} are unchanged by it.

\subsection*{v4.8.0-r07 working assignment}
\label{subsec:v480-r07-assignment}

For \(z\in\mathbb C^3\) let \(s(z)=z_1^2+z_2^2+z_3^2\) as in r06.  The working assignment is
\[
s>0\ \ \text{matter},\qquad s<0\ \ \text{tachyon},\qquad s=0\ \ \text{photon}.
\]
It is adopted as the author's working choice (2026-10-19) and can be reversed without affecting
the theoremlets below, which are stated for the sign of \(s\).

\paragraph{Motivation.}
Layer r82 types the tachyon frame by the signed scale \(\ell_T^2=-\lambda_T^2\), \(\lambda_T>0\),
an allowed effective statement of the release scope.  Under the working assignment the sign of
\(s\) carries exactly this interval signature: positive for matter, negative for tachyon, zero for
the photon.  It is not the orientation code of
\S\ref{subsec:v480-r01-orientation-vs-signature}, which r01 keeps separate from the r82 signature.

\paragraph{Consequence for the earlier classification.}
The classification of \S\ref{sec:unitary} placed the photon kernel at purely imaginary \(z\).  In
Reading~B a purely imaginary non-zero \(z=iy\) has \(s=-|y|^2<0\) and is tachyon-type, not
light-like.  The earlier classification is therefore replaced, not translated.

\subsection*{v4.8.0-r07 duality angle}
\label{subsec:v480-r07-duality-angle}

\paragraph{Theoremlet v4.8.0-r07.1.}
Every \(z\in\mathbb C^3\) with \(s(z)\neq0\) can be written as
\[
z=e^{i\alpha}z_m,\qquad s(z_m)=|s(z)|>0,
\]
with \(\alpha=\tfrac12\arg s(z)\) determined modulo \(\pi\) and \(z_m\) determined up to sign.
The configuration is of matter type exactly when \(\alpha\equiv0\) and of tachyon type exactly when
\(\alpha\equiv\pi/2\pmod\pi\).

\paragraph{Proof sketch.}
\(s(e^{i\alpha}z_m)=e^{2i\alpha}s(z_m)\), so \(s(z_m)=e^{-2i\alpha}s(z)\) is real positive iff
\(2\alpha\equiv\arg s(z)\pmod{2\pi}\).  Replacing \(\alpha\) by \(\alpha+\pi\) replaces \(z_m\) by
\(-z_m\).  The two special values follow from the sign of \(s\).  Checked numerically for 500 random
configurations.  \(\square\)

\paragraph{Theoremlet v4.8.0-r07.2.}
Under the group \(SO(3,\mathbb C)\times U(1)_{\rm dual}\), acting by \(z\mapsto e^{i\alpha}Rz\), the
orbits in \(\mathbb C^3\setminus\{0\}\) are the shells \(\{|s|=r\}\), \(r>0\), and the punctured null
cone.  In particular, a matter configuration and a tachyon configuration with the same \(|s|\) lie in
one orbit, related by \(\alpha=\pi/2\); the null cone is the only orbit fixed by the whole duality
circle.

\paragraph{Proof sketch.}
By Theoremlet r06.3 the \(SO(3,\mathbb C)\)-orbits are the level sets of \(s\); the duality circle
rotates \(s\) through all values of fixed modulus and fixes \(s=0\).  Transitivity on each shell was
checked constructively for 200 random pairs.  \(\square\)

\paragraph{Reading (conjecture).}
Non-real \(s\) is not a further sector: it is a matter configuration seen in a rotated duality frame.
In this precise sense matter and tachyon are one object in two duality frames, distinguished
Lorentz-invariantly but exchanged by duality, while the photon is the duality-invariant fix-sector.
This is the form the initial conjecture of Theorem~\ref{thm:unity} takes in Reading~B.  Whether the
duality angle is physical, or a gauge freedom to be fixed, is not decided here.

\subsection*{v4.8.0-r07 the parameter \texorpdfstring{$t$}{t}}
\label{subsec:v480-r07-parameter-t}

\paragraph{Observation.}
On the null set of the natural norm, \(N(q)=-t^2+s(z)=0\), one has \(s=t^2\).  For real
\(t\neq0\) this forces \(\alpha\equiv0\), the matter frame, with \(t=\sqrt{|s|}\); for \(t=0\) it forces
the photon sector; a tachyon-type kernel on this set requires \(t^2<0\).  The condition \(N(q)=0\)
therefore ties \(t\) to the invariant magnitude \(\sqrt{|s|}\) and fixes the duality frame.  Whether
the Sphaera kernel is meant to satisfy \(N(q)=0\) is open; the observation is recorded as a candidate
role for \(t\), not adopted.

\subsection*{v4.8.0-r07 ledger and guard update}
\label{subsec:v480-r07-ledger-guard}

\[
R_{\rm sect}^{v480\text{-}r07}=(0,0,0,1,1)
\]
for (working assignment adopted; duality angle; orbit structure; physical status of the duality angle
open; role of \(t\) open).

The layer spends no guard:
\[
R_{\rm glob}^{(v480\text{-}r07,+)}=(0,1,1,1,1),
\qquad
\Gamma_{\rm guard}^{v480\text{-}r07}=(1,1,1,1,1),
\qquad
\beta_4=1/256\ \text{unspent}.
\]
It does not imply exact finite \(C5\) zero closure, exact all-readout terminality,
\(C6/T2\), a global photon-clock/RH theorem, or global all-readout closure.


% ============================================================================
\clearpage
\section*{v4.8.0-r08 Source-Only: Test of Prediction P1 (Photon-Frame Unfolding)}
\addcontentsline{toc}{section}{v4.8.0-r08 Source-Only: Test of Prediction P1 (Photon-Frame Unfolding)}
\label{sec:v480-r08-p1-test}
% ============================================================================

\paragraph{Purpose.}
Corollary~\ref{cor:falsifiable} lists P1 as testable now: unfolding the zeta zeros with the
photon-frame density should improve the GUE fit by about \(21\%\) (Corollary~\ref{cor:odlyzko}),
and Theorem~\ref{thm:gue-selfdet} states that the Kolmogorov--Smirnov distance has a unique global
minimum at \(\mu=\gL\).  This layer performs the test on the first 2000 zeros.  Earlier layers are
not rewritten.

\noindent\textbf{Source-only status.}
This v4.8.0-r08 layer is source-only.  It does not build a PDF, does not assign a new DOI,
and does not authorize a public release.  The v4.8.0 Version DOI
\href{https://doi.org/10.5281/zenodo.23365244}{10.5281/zenodo.23365244} and the Concept DOI
\href{https://doi.org/10.5281/zenodo.19210728}{10.5281/zenodo.19210728} are unchanged by it.

\subsection*{v4.8.0-r08 setup}
\label{subsec:v480-r08-setup}

As in Theorem~\ref{thm:gue-selfdet},
\(\tilde\gamma_n(\mu)=\int_0^{\gamma_n}\frac{\log(t/2\pi\mu)}{2\pi\mu}\,dt
=\frac{\gamma_n}{2\pi\mu}\bigl(\log\frac{\gamma_n}{2\pi\mu}-1\bigr)\) and
\(s_n(\mu)=\tilde\gamma_{n+1}(\mu)-\tilde\gamma_n(\mu)\).  The comparison law is the Wigner surmise
\(p(s)=\frac{32}{\pi^2}s^2e^{-4s^2/\pi}\); \(\chi^2\)/dof uses 15 bins on \([0,3]\).  The scale runs over
\(\mu\in[0.5,3]\) in steps of \(0.05\).  Two readings are tested: (a) raw, exactly as defined, and
(b) normalized to mean spacing \(1\), since a pure rescaling would otherwise move the mean.

\subsection*{v4.8.0-r08 result}
\label{subsec:v480-r08-result}

\begin{center}
\small
\begin{tabular}{@{}>{\raggedright\arraybackslash}p{0.26\textwidth}>{\raggedright\arraybackslash}p{0.33\textwidth}>{\raggedright\arraybackslash}p{0.33\textwidth}@{}}
\toprule
 & zeros 1--100 & zeros 101--2000\\
\midrule
(a) raw: KS minimum & \(\mu=1.00\) & \(\mu=1.00\)\\
(a) KS at \(\mu=1\) / \(\mu=\gL\) & \(0.068\) / \(0.699\) & \(0.041\) / \(0.628\)\\
(a) mean spacing at \(\mu=\gL\) & \(0.419\) & \(0.473\)\\
(b) normalized: KS minimum & \(\mu=2.85\) & \(\mu=2.25\)\\
(b) KS range over \(\mu\) & \(0.037\)--\(0.070\) & \(0.039\)--\(0.042\)\\
(b) \(\chi^2\)/dof change \(\mu=1\to\gL\) & \(+5.4\%\) improvement & \(-9.9\%\) (worse)\\
\bottomrule
\end{tabular}
\end{center}

\paragraph{Theoremlet v4.8.0-r08.1.}
On the first 2000 zeros, neither reading of the photon-frame unfolding singles out \(\mu=\gL\).
In reading (a) the unique minimum is the standard unfolding \(\mu=1\); in reading (b) the
dependence on \(\mu\) is flat within the statistical scatter and its minimum is not at \(\gL\).
Moreover, in reading (b) the effect of \(\mu\) vanishes with height: the local stretch it introduces
is of relative size \(\log\mu/\log(\gamma/2\pi)\), about \(1\%\) at \(\gamma\sim10^{22}\), so an
asymptotic improvement of \(21\%\) cannot occur.

\paragraph{Proof sketch.}
Direct computation, script \texttt{k11\_p1\_entfaltung.py} in the working folder.  For reading (b),
\(\rho_\mu(t)/\rho_1(t)=\mu^{-1}\bigl(1-\log\mu/\log(t/2\pi)\bigr)\); after normalization the constant
factor \(\mu^{-1}\) drops out and only the second term, which decays like \(1/\log t\), remains.
\(\square\)

\paragraph{Reproducibility note.}
The value of \(75.8\%\) for 100 zeros stated in Corollary~\ref{cor:odlyzko} is not reproduced in
either reading.  The original script \texttt{gue\_frame\_correction.py} is not in the source folders,
so its exact procedure could not be compared.

\subsection*{v4.8.0-r08 consequences}
\label{subsec:v480-r08-consequences}

\begin{center}
\small
\begin{tabular}{@{}>{\raggedright\arraybackslash}p{0.34\textwidth}>{\raggedright\arraybackslash}p{0.58\textwidth}@{}}
\toprule
statement & status after r08\\
\midrule
Thm.~\ref{thm:gue-selfdet} (unique minimum at \(\gL\)) & not confirmed; contradicted on zeros 1--2000\\
Cor.~\ref{cor:odlyzko}, P1 in Cor.~\ref{cor:falsifiable} & by its own criterion falsified at low height; the asymptotic form excluded by the \(1/\log\gamma\) decay\\
\(\gL\) as ``GUE fixpoint of the frame'' & withdrawn; \(\gL\) remains an input (cf.\ r03 F11, OP1)\\
\bottomrule
\end{tabular}
\end{center}

Corollary~\ref{cor:odlyzko} states that non-observation falsifies \(\gL=\gBI/\gstring\) as the
physical frame factor.  This layer records the non-observation on the accessible zeros and the
reason why the stated asymptotic effect cannot appear.  It does not by itself decide the physical
status of \(\gL\); that remains with OP1.

\subsection*{v4.8.0-r08 ledger and guard update}
\label{subsec:v480-r08-ledger-guard}

\[
R_{\rm P1}^{v480\text{-}r08}=(0,0,1)
\]
for (test performed; prediction not confirmed; physical status of \(\gL\) open).

The layer spends no guard:
\[
R_{\rm glob}^{(v480\text{-}r08,+)}=(0,1,1,1,1),
\qquad
\Gamma_{\rm guard}^{v480\text{-}r08}=(1,1,1,1,1),
\qquad
\beta_4=1/256\ \text{unspent}.
\]
It does not imply exact finite \(C5\) zero closure, exact all-readout terminality,
\(C6/T2\), a global photon-clock/RH theorem, or global all-readout closure.


% ============================================================================
\clearpage
\section*{v4.8.0-r09 Source-Only: Tests of P2 and P3, and Re-Classification of Earlier Results}
\addcontentsline{toc}{section}{v4.8.0-r09 Source-Only: Tests of P2 and P3, and Re-Classification of Earlier Results}
\label{sec:v480-r09-p2-p3-reclassification}
% ============================================================================

\paragraph{Purpose.}
Corollary~\ref{cor:falsifiable} lists P2 (Koide phase) as ``already matches data'' and P3
(\(\bsym\)-formula) as ``matches to \(0.018\%\)''.  This layer tests both against the measurement
precision and against the freedom in choosing the formula.  Since layers r03--r08 changed the status
of many earlier statements, it then collects the current status of all predictions and of the main
results in one register, so that no earlier ``proved'' or ``\checkmark'' is read without its
present qualification.  Earlier layers are not rewritten.

\noindent\textbf{Source-only status.}
This v4.8.0-r09 layer is source-only.  It does not build a PDF, does not assign a new DOI,
and does not authorize a public release.  The v4.8.0 Version DOI
\href{https://doi.org/10.5281/zenodo.23365244}{10.5281/zenodo.23365244} and the Concept DOI
\href{https://doi.org/10.5281/zenodo.19210728}{10.5281/zenodo.19210728} are unchanged by it.

\subsection*{v4.8.0-r09 P2: Koide phase}
\label{subsec:v480-r09-p2}

Theorem~\ref{thm:koide-triolic} predicts, in the Brannen form
\(\sqrt{m_k}=A\bigl(1+\sqrt2\cos(\theta_K+2\pi k/3)\bigr)\), the phase
\(\theta_K=\pi-\sqrt2\arcsin\bsym=2.319109\) rad \(=132.8751^\circ\).

\paragraph{Theoremlet v4.8.0-r09.1.}
From \(m_e=0.51099895\), \(m_\mu=105.6583755\) and \(m_\tau=1776.93\pm0.09\) MeV the observed phase is
\(\theta_K^{\rm obs}=132.7325^\circ\pm0.0004^\circ\) (with \(m_\tau=1776.86\pm0.12\):
\(132.7328^\circ\pm0.0005^\circ\)).  The prediction differs by \(+0.142^\circ\), i.e.\ by roughly
\(300\)--\(400\) standard deviations.  P2 is excluded by the data.

\paragraph{Proof sketch.}
Solve the Brannen form for \(\theta\) from the three masses; propagate the \(\tau\)-mass error, which
dominates.  Script \texttt{k12\_p2\_p3.py} in the working folder.  \(\square\)

\paragraph{Remarks.}
(i) The relative deviation of \(0.11\%\) quoted in Theorem~\ref{thm:koide-triolic} is small compared with
the angle, but large compared with the measurement error; a prediction is tested against the error.
(ii) The observed phase is \(\theta_K^{\rm obs}=2\pi/3+0.22223\), within \(10^{-5}\) of
\(2\pi/3+2/9\); this empirical regularity is known in the literature and is not produced by the formula
of Theorem~\ref{thm:koide-triolic}.
(iii) Step~1 of the proof of Theorem~\ref{thm:koide-triolic} derives \(K=2/3\) from Cauchy--Schwarz.
Cauchy--Schwarz gives \((\sum\sqrt{m_i})^2\le3\sum m_i\), i.e.\ \(K\ge1/3\), with equality only for equal
masses; it does not yield \(K=2/3\).  The Koide value \(K=0.666665\) is an empirical fact; its derivation
is open.

\subsection*{v4.8.0-r09 P3: the \texorpdfstring{$\bsym$}{beta}-formula}
\label{subsec:v480-r09-p3}

Remark~\ref{conj:cl-correction} states \(\bsym=\tfrac1{\sqrt3}\exp\bigl(-\gamma_1/(4\pi\gL^5)\bigr)\).

\paragraph{Theoremlet v4.8.0-r09.2.}
With \(\gL=\gBI/\gstring=1.86444\) the formula gives \(0.549233\) against
\(\bsym=\sqrt{(\gL-1)/(\gL+1)}=0.549348\), a deviation of \(-0.021\%\); it holds exactly for
\(\gL=1.86366\).  It is a relation between the two inputs \(\gBI,\gstring\) and \(\gamma_1\), not a
prediction of an independently measured quantity.  Among the 56 formulas
\(\tfrac1{\sqrt3}\exp(-\gamma_1/(c\,\gL^n))\), \(n=1,\dots,8\),
\(c\in\{1,2,\pi,2\pi,4\pi,8\pi,16\pi\}\), 27 lie within \(\pm5\%\); the chance that one of these lands
within \(\pm0.02\%\) is about \(11\%\).

\paragraph{Proof sketch.}
Both sides are functions of \(\gL\) alone; solve by bisection.  Enumerate the family and estimate the
look-elsewhere probability with a uniform density over \(\pm5\%\).  \(\square\)

\paragraph{Assessment.}
The exponent was changed from \(\gL^2\) to \(\gL^5\) after the fact (Remark~\ref{conj:cl-correction}),
and Remark~\ref{conj:cl-correction} itself marks it as requiring structural derivation.  P3 is
recorded as a numerical coincidence of moderate significance, not as a confirmed prediction.  A
derivation of the exponent before comparison with data would change this.

\subsection*{v4.8.0-r09 register of predictions}
\label{subsec:v480-r09-predictions}

\begin{center}
\small
\begin{tabular}{@{}>{\raggedright\arraybackslash}p{0.27\textwidth}>{\raggedright\arraybackslash}p{0.31\textwidth}>{\raggedright\arraybackslash}p{0.34\textwidth}@{}}
\toprule
prediction (Cor.~\ref{cor:falsifiable}) & earlier status & status after r09\\
\midrule
P1 GUE photon-frame unfolding & testable now & not confirmed (r08); minimum at \(\mu=1\)\\
P2 Koide phase & matches to \(0.111\%\) & excluded by data, \(\sim300\)--\(400\sigma\) (r09.1)\\
P3 \(\bsym\)-formula & matches to \(0.018\%\) & coincidence, exponent chosen post hoc; \(\sim11\%\) look-elsewhere (r09.2)\\
P4 \(\Lambda\) as frame artefact & reinterpretation & contradicted as stated (r03 F4); \(\Lambda\) type open (r03 F3); \(\Theta\)-sum route false (r05)\\
P5 \(\alpha\)-drift ratio & not yet testable & exponent corrected (r03 F2); untestable at present\\
P6 Hawking, tachyon frame & theoretical & assumption \(k(M)=0\) made explicit (r03 F6); untested\\
\bottomrule
\end{tabular}
\end{center}

\subsection*{v4.8.0-r09 register of main results}
\label{subsec:v480-r09-results}

\begin{center}
\small
\begin{tabular}{@{}>{\raggedright\arraybackslash}p{0.34\textwidth}>{\raggedright\arraybackslash}p{0.58\textwidth}@{}}
\toprule
statement & current status (layer)\\
\midrule
Thm.~\ref{thm:unity}, equivalence of (i)--(v) & conjecture; Stone--von Neumann justification withdrawn (r04); Reading~B gives a precise form via duality (r07)\\
Thm.~\ref{thm:unitary-equiv} & open (r04)\\
Def.~\ref{def:biquaternion}, Minkowski reading & replaced: signature \((1,6)\); natural norm \(N(q)\), Reading~B (r06)\\
Thm.~\ref{thm:lorentz} & shows \(J\)-invariance; Lorentz action from Reading~B (r06)\\
Thm.~\ref{thm:hopf} & Hopf fibrations classical; physical identification conjecture; the chain \(S^{15}\to S^8\to S^4\to S^2\) is not a composition of Hopf maps (recorded here)\\
Thm.~\ref{thm:fractal} & null axis redefined as \(s=0\), directions \(=S_0\) (r06); transport open\\
triole structure & triole lemma proved (r06.2); octonionic triole configuration repaired (r04.4)\\
Heisenberg compensator, cascade & standalone, proved (r04.2, r04.3)\\
Thm.~\ref{thm:einstein} & structural conjecture; action undefined (r05)\\
Thm.~\ref{thm:gue} & empirical fact; derivation withdrawn (r05)\\
Thm.~\ref{thm:gue-selfdet} & not confirmed (r08)\\
Thm.~\ref{thm:koide-triolic} & \(K=2/3\) empirical, derivation open; phase formula excluded (r09)\\
Thm.~\ref{thm:lambda-zero} & (i) holds by definition of the photon frame; (iii) value off by \(10^{121}\) (r03 F8, F9)\\
Cor.~\ref{cor:matter-error} & corrected to \(\gL^5\approx22.5\) (r03 F1)\\
\(\gL=\gBI/\gstring\) & ratio of two published Barbero--Immirzi values (r03 F11); input, not derived (OP1)\\
OP5 Step~(B) & no valid form (r05)\\
RH & not proved; forbidden as global claim (\(W_\gamma^{\rm glob}\)); Riemann source no proof source (r04)\\
\bottomrule
\end{tabular}
\end{center}

\paragraph{Theoremlet v4.8.0-r09.3.}
After r09 no prediction of Corollary~\ref{cor:falsifiable} stands as confirmed.  The results that carry
are the structural ones: the standalone compensator and cascade, the repaired triole hierarchy, the
triole lemma, the signature analysis and Reading~B with its orbit, null-cone and duality structure.
None of the five guards is affected.

\paragraph{Proof sketch.}
Collect the decisions of r03--r09 as listed in the two registers.  \(\square\)

\subsection*{v4.8.0-r09 ledger and guard update}
\label{subsec:v480-r09-ledger-guard}

\[
R_{\rm pred}^{v480\text{-}r09}=(0,0,0,1)
\]
for (P2 decided; P3 classified; register compiled; derivations of \(K=2/3\) and of the P3 exponent open).

The layer spends no guard:
\[
R_{\rm glob}^{(v480\text{-}r09,+)}=(0,1,1,1,1),
\qquad
\Gamma_{\rm guard}^{v480\text{-}r09}=(1,1,1,1,1),
\qquad
\beta_4=1/256\ \text{unspent}.
\]
It does not imply exact finite \(C5\) zero closure, exact all-readout terminality,
\(C6/T2\), a global photon-clock/RH theorem, or global all-readout closure.


% ============================================================================
\clearpage
\section*{v4.8.0-r10 Source-Only: The Koide Structure in the Complex-Vector Reading}
\addcontentsline{toc}{section}{v4.8.0-r10 Source-Only: The Koide Structure in the Complex-Vector Reading}
\label{sec:v480-r10-koide-reading-b}
% ============================================================================

\paragraph{Purpose.}
Layer r09 excluded the phase formula of Theorem~\ref{thm:koide-triolic} and noted that the observed
Koide phase lies close to \(2\pi/3+2/9\).  This layer asks what Reading~B (r06, r07) says about the
lepton masses before any number is fitted.  Because the observed values are known, nothing found here
can count as a prediction; the layer records what the structure fixes and what it leaves free.
Earlier layers are not rewritten.

\noindent\textbf{Source-only status.}
This v4.8.0-r10 layer is source-only.  It does not build a PDF, does not assign a new DOI,
and does not authorize a public release.  The v4.8.0 Version DOI
\href{https://doi.org/10.5281/zenodo.23365244}{10.5281/zenodo.23365244} and the Concept DOI
\href{https://doi.org/10.5281/zenodo.19210728}{10.5281/zenodo.19210728} are unchanged by it.

\subsection*{v4.8.0-r10 the Koide form is a triole}
\label{subsec:v480-r10-koide-triole}

\paragraph{Theoremlet v4.8.0-r10.1.}
For \(v=(\sqrt{m_e},\sqrt{m_\mu},\sqrt{m_\tau})\) the Brannen form
\(\sqrt{m_k}=A\bigl(1+\sqrt2\cos(\theta+2\pi k/3)\bigr)\) is identical to
\[
v=t\,(1,1,1)+\sqrt2\,\operatorname{Re}\bigl(c\,(1,\omega,\omega^2)\bigr),\qquad
t=A,\quad c=|c|e^{i\theta},
\]
a scalar plus the real part of a symmetric triole.  For any \(t,c\),
\(K=\sum m_k/(\sum\sqrt{m_k})^2=\tfrac13\bigl(1+|c|^2/t^2\bigr)\); hence \(K=2/3\) if and only if
\(|c|=t\).  From the measured masses, \(|c|/t=0.999997\) and \(\theta\bmod 2\pi/3=0.222225\).

\paragraph{Proof sketch.}
Discrete Fourier decomposition of \(v\) over \(\mathbb Z/3\): the constant mode is \(t\), the first mode
is \(c\); the reconstruction is exact.  \(\sum v_k^2=3t^2+3|c|^2\) and \(\sum v_k=3t\) give \(K\).
\(\square\)

\paragraph{Theoremlet v4.8.0-r10.2.}
Embed the triole as the kernel part of Reading~B, \(q=it\cdot\mathbf 1+\sum_k z_ke_k\) with
\(z=c(1,\omega,\omega^2)\).  Then \(s(z)=0\) and \(N(q)=-t^2\) for every \(c\); the Reading-B invariants
therefore fix neither \(|c|/t\) (i.e.\ \(K\)) nor \(\arg c\) (the Koide phase).

\paragraph{Proof sketch.}
\(s=c^2(1+\omega^2+\omega^4)=0\) by the triole lemma r06.2, and \(N(q)=-t^2+s\).  \(\square\)

\paragraph{Observation (conjecture).}
The condition \(|c|=t\) is equivalent to \(-t^2+|z_k|^2=0\) for each kernel separately, i.e.\ each
kernel is light-like in the original per-kernel norm of Definition~\ref{def:biquaternion}.  This
equivalence depends on the factor \(\sqrt2\) of the Brannen form, which is a convention; it is
recorded as a candidate principle for \(K=2/3\), not as a derivation.

\subsection*{v4.8.0-r10 no canonical angle gives the phase}
\label{subsec:v480-r10-phase}

The following angles of the framework were fixed before comparison: \(j\pi/12\) for \(j=0,\dots,7\);
\(\arcsin\bsym\), \(2\arcsin\bsym\), \(\arcsin(1/\sqrt3)\), \(2\arcsin(1/\sqrt3)\), \(\arctan\bsym\),
\(\bsym\) and \(1/\sqrt3\) read as radians; the P2 angle; the tetrahedral and magic angles.  Reduced
modulo \(2\pi/3\), the closest is the excluded P2 value (\(0.2247\), \(356\sigma\) away), the next
\(\pi/12\) (\(0.2618\)).  None reproduces \(\delta_{\rm obs}=0.222225\).  The value \(2/9\) is not a
rational multiple of \(\pi\), which is what angular structures of this type produce.

\subsection*{v4.8.0-r10 what is missing: dynamics}
\label{subsec:v480-r10-dynamics}

Reading~B is kinematic: it fixes how the kernels transform, not which configuration is realised.

\paragraph{Theoremlet v4.8.0-r10.3.}
The quadratic Lorentz-invariant real functionals of \(z\in\mathbb C^3\) are
\(\operatorname{Re}(e^{-i\varphi}s)=\cos\varphi\,(E^2-B^2)+\sin\varphi\,2E\cdot B\) for \(z=E+iB\), and a
duality rotation \(z\mapsto e^{ia}z\) maps \(\varphi\mapsto\varphi-2a\).

\paragraph{Proof sketch.}
\(s\) is the only quadratic \(SO(3,\mathbb C)\)-invariant up to complex scale; its real-linear functionals
are the real parts of \(e^{-i\varphi}s\).  \(s(e^{ia}z)=e^{2ia}s(z)\).  Checked numerically.  \(\square\)

\paragraph{Reading.}
\(\varphi=0\) is the Maxwell Lagrangian; \(\sin\varphi\,2E\cdot B\) is the topological \(\theta\)-type term.
The simplest dynamics of Reading~B is thus electromagnetism, consistent with the photon sector, while the
duality angle of r07 reappears as a vacuum angle.  Lepton masses, \(K=2/3\) and the Koide phase require an
additional principle that fixes \(|c|/t\) and \(\arg c\).  Any such principle has to be stated before it is
compared with the known values; only then can it yield a prediction.

\subsection*{v4.8.0-r10 ledger and guard update}
\label{subsec:v480-r10-ledger-guard}

\[
R_{\rm Koide}^{v480\text{-}r10}=(0,0,0,1,1)
\]
for (Koide form identified as triole; Reading-B invariants shown not to fix it; canonical angles excluded;
principle for \(|c|=t\) open; principle for the phase open).

The layer spends no guard:
\[
R_{\rm glob}^{(v480\text{-}r10,+)}=(0,1,1,1,1),
\qquad
\Gamma_{\rm guard}^{v480\text{-}r10}=(1,1,1,1,1),
\qquad
\beta_4=1/256\ \text{unspent}.
\]
It does not imply exact finite \(C5\) zero closure, exact all-readout terminality,
\(C6/T2\), a global photon-clock/RH theorem, or global all-readout closure.


% ============================================================================
\clearpage
\section*{v4.8.0-r11 Source-Only: Pre-Registered Test of the Phase Relation \texorpdfstring{$\delta=K/3$}{delta = K/3}}
\addcontentsline{toc}{section}{v4.8.0-r11 Source-Only: Pre-Registered Test of the Phase Relation delta = K/3}
\label{sec:v480-r11-preregistered-h1}
% ============================================================================

\paragraph{Purpose.}
Layer r10 wrote the lepton masses as a scalar plus a triole and left the phase open.  For the charged
leptons the phase satisfies \(\delta\approx K/3\approx2/9\).  This layer records the first
pre-registered test of the Companion: the hypothesis that \(\delta=K/3\) holds for every triplet of
charged fermions.  Hypothesis, definitions, data and decision criterion were committed to the
repository before any computation (working-folder note \texttt{Vorab\_Notiz\_H1\_Koide\_Phase.md},
commit \texttt{81f38d3}).  Earlier layers are not rewritten.

\noindent\textbf{Source-only status.}
This v4.8.0-r11 layer is source-only.  It does not build a PDF, does not assign a new DOI,
and does not authorize a public release.  The v4.8.0 Version DOI
\href{https://doi.org/10.5281/zenodo.23365244}{10.5281/zenodo.23365244} and the Concept DOI
\href{https://doi.org/10.5281/zenodo.19210728}{10.5281/zenodo.19210728} are unchanged by it.

\subsection*{v4.8.0-r11 hypothesis and protocol}
\label{subsec:v480-r11-protocol}

\paragraph{Hypotheses.}
H1: \(\delta=K/3\) for each triplet.  H1\('\): \(\delta=2/9\) for each triplet.  For quarks
\(K\neq2/3\), so the two differ.  The factor \(1/3\) may be read as \(1/N\) for \(N=3\) kernels or as
\(\beta_{\rm sym,Hopf}^2\); for triplets both readings predict the same number.

\paragraph{Protocol.}
\(v_k=\sqrt{m_k}\) in ascending order, all roots positive; \(t\) and \(c\) as in Theoremlet r10.1;
\(\delta=\arg c\) reduced to \([0,2\pi/3)\).  Primary data: running \(\overline{\rm MS}\) Yukawa couplings
at \(\mu=M_Z\) with PDG-2024 inputs (Antusch, Hinze, Saad, arXiv:2510.01312), for which masses are
proportional to couplings and the common factor cancels.  Error propagation by \(10^6\) Monte-Carlo
draws.  Criterion: \(|z|<2\) compatible, \(|z|>3\) refuted; H1 is refuted if it is refuted for at least
one quark triplet.

\subsection*{v4.8.0-r11 result}
\label{subsec:v480-r11-result}

\begin{center}
\small
\begin{tabular}{@{}>{\raggedright\arraybackslash}p{0.27\textwidth}rrrrr@{}}
\toprule
triplet (at \(M_Z\)) & \(K\) & \(\delta\) & \(K/3\) & \(z\)(H1) & \(z\)(H1\('\))\\
\midrule
up quarks \((u,c,t)\) & 0.8876 & 0.0518 & 0.2959 & \(-322\) & \(-353\)\\
down quarks \((d,s,b)\) & 0.7478 & 0.1001 & 0.2493 & \(-115\) & \(-128\)\\
leptons \((e,\mu,\tau)\) & 0.6678 & 0.2211 & 0.2226 & \(-49\) & \(-46\)\\
leptons, pole masses & 0.6667 & 0.2222 & 0.2222 & \(+0.4\) & \(+0.4\)\\
\bottomrule
\end{tabular}
\end{center}

\paragraph{Theoremlet v4.8.0-r11.1.}
By the pre-registered criterion both H1 and H1\('\) are refuted.  The relation
\(\delta\approx K/3\approx2/9\) holds only for the charged leptons with pole masses, where it served as
calibration; it fails for both quark triplets and already for the running lepton masses at \(M_Z\).

\paragraph{Proof sketch.}
Direct computation (script \texttt{k14\_test\_h1.py}, seed 20261020), following the committed protocol
without deviation.  The secondary data set (PDG-2024 masses at their standard scales) gives the same
verdict.  \(\square\)

\paragraph{Assessment.}
The lepton coincidence is tied to pole masses, as is the Koide relation itself, and is not a law of the
triole structure for all triplets.  Hypotheses read off from the quark values found here would again be
post-dictions and would need their own pre-registration.  The procedure itself is retained: it is the
Companion's first test that could have failed and was fixed before the data were examined.

\subsection*{v4.8.0-r11 ledger and guard update}
\label{subsec:v480-r11-ledger-guard}

\[
R_{\rm H1}^{v480\text{-}r11}=(0,0,1)
\]
for (protocol committed before computation; hypothesis decided; principle for the triole parameters
still open).

The layer spends no guard:
\[
R_{\rm glob}^{(v480\text{-}r11,+)}=(0,1,1,1,1),
\qquad
\Gamma_{\rm guard}^{v480\text{-}r11}=(1,1,1,1,1),
\qquad
\beta_4=1/256\ \text{unspent}.
\]
It does not imply exact finite \(C5\) zero closure, exact all-readout terminality,
\(C6/T2\), a global photon-clock/RH theorem, or global all-readout closure.


% ============================================================================
\clearpage
\section*{v4.8.0-r12 Source-Only: Pre-Registered Neutrino Tests, the Octonionic Flavour Group, and Reclassification of the PMNS Claims}
\addcontentsline{toc}{section}{v4.8.0-r12 Source-Only: Pre-Registered Neutrino Tests, Octonionic Flavour Group, PMNS Reclassification}
\label{sec:v480-r12-neutrino-flavour}
% ============================================================================

\paragraph{Purpose.}
Layer r11 introduced pre-registration.  This layer records the seven further pre-registered tests
H2--H8, which apply the triole form to neutrinos and to lepton mixing, the structural facts they rest on,
and a reclassification of the earlier PMNS statements of the Companion.  For every test the hypothesis,
data and decision criterion were committed to the repository before the computation; the commits are
listed below.  Earlier layers are not rewritten.

\noindent\textbf{Source-only status.}
This v4.8.0-r12 layer is source-only.  It does not build a PDF, does not assign a new DOI,
and does not authorize a public release.  The v4.8.0 Version DOI
\href{https://doi.org/10.5281/zenodo.23365244}{10.5281/zenodo.23365244} and the Concept DOI
\href{https://doi.org/10.5281/zenodo.19210728}{10.5281/zenodo.19210728} are unchanged by it.

\subsection*{v4.8.0-r12 structural facts}
\label{subsec:v480-r12-structure}

\paragraph{Theoremlet v4.8.0-r12.1 (triole spectrum and sign).}
For \(v_k=t\bigl(1+\sqrt2\,r\cos(\delta+2\pi k/3)\bigr)\) one has
\(v_0v_1v_2=t^3\bigl(1-\tfrac32r^2+\tfrac{r^3}{\sqrt2}\cos3\delta\bigr)\).  For \(r=1\) at most one
component is negative, and exactly one is negative if and only if \(\cos3\delta<1/\sqrt2\); the boundary
\(\delta=\pi/12\) is the point where one mass vanishes.  The vector \(v\) is the spectrum of the Hermitian
circulant matrix \(t\,\mathbf 1+\tfrac{c}{\sqrt2}P+\tfrac{\bar c}{\sqrt2}P^2\) (\(P\) the cyclic shift),
whose eigenvectors are the fixed Fourier vectors.

\paragraph{Proof sketch.}
Expand the product with \(\sum_k\cos\theta_k=0\), \(\sum_k\cos^2\theta_k=3/2\) and
\(\prod_k\cos\theta_k=\tfrac14\cos3\delta\) for \(\theta_k=\delta+2\pi k/3\).  The circulant statement is the
discrete Fourier diagonalization.  The identities were also checked numerically to \(10^{-15}\).  \(\square\)

\paragraph{Theoremlet v4.8.0-r12.2 (octonionic flavour group).}
In the Cayley--Dickson convention of layer r04 the automorphisms of \(\OO\) that permute
\(e_1,\dots,e_7\) up to sign form a group of order \(1344\).  Its stabilizer of \(+e_7\) has order \(96\),
acts unitarily with determinant one and irreducibly on \(\mathbb C^3=\mathrm{span}_{\mathbb C}(e_4,e_5,e_6)\)
with complex structure \(L_{e_7}\), and has the same element-order statistics and the same number of
conjugacy classes (ten) as \(\Delta(96)\subset SU(3)\).

\paragraph{Proof sketch.}
An automorphism is fixed by the images of \(e_1,e_2,e_4\); all \(14\cdot12\cdot8\) choices extend and
preserve the full multiplication table.  Restriction, unitarity, determinant, irreducibility
(\(\langle|\chi|^2\rangle=1\)) and the class data were computed exhaustively (scripts
\texttt{k22\_oktonion\_triplett.py}, \texttt{k23\_test\_h6.py}).  The identification with \(\Delta(96)\) is by
these invariants; an explicit isomorphism is not written out here.  \(\square\)

\paragraph{Theoremlet v4.8.0-r12.3 (residual reflections).}
With the triole \(e_4\to e_5\to e_6\) as charged-lepton symmetry and the electron along its invariant
direction, the fifteen involutions of the stabilizer preserve the following columns of the lepton mixing
matrix:
\begin{center}
\small
\begin{tabular}{>{\raggedright\arraybackslash}p{0.40\textwidth}>{\raggedright\arraybackslash}p{0.07\textwidth}>{\raggedright\arraybackslash}p{0.20\textwidth}>{\raggedright\arraybackslash}p{0.21\textwidth}}
\toprule
Involution & count & preserved \(|U_{\alpha i}|^2\) & pattern\\
\midrule
pure sign changes (one generation fixed) & 3 & \((\tfrac13,\tfrac13,\tfrac13)\) & TM2\\
Fano centre in \(\{e_1,e_2,e_3\}\), first lift & 3 & \((\tfrac23,\tfrac16,\tfrac16)\) & TM1\\
Fano centre in \(\{e_1,e_2,e_3\}\), second lift & 3 & \((0,\tfrac12,\tfrac12)\) & \(\theta_{13}=0\)\\
Fano centre in \(\{e_4,e_5,e_6\}\) & 6 & \((\tfrac13,\tfrac{2\mp\sqrt3}{6})\) & \(\Delta(96)\) column\\
Fano centre \(e_7\) & 0 & --- & only order-4 lifts\\
\bottomrule
\end{tabular}
\end{center}
The TM1 lift and the \(\theta_{13}=0\) lift over the same Fano image differ by a pure sign change.

\paragraph{Theoremlet v4.8.0-r12.4 (Cayley--Dickson involution as CP).}
The involution \(p+q\ell\mapsto p-q\ell\) commutes with every lift of the triole and with all TM1 and
\(\theta_{13}=0\) involutions, and acts on the triplet as \(z\mapsto-\bar z\) in the basis in which the
triole is the plain cyclic shift.  If the Majorana mass matrix is invariant under it and under a residual
reflection, then \(\sin^2\theta_{23}=\tfrac12\) and \(\delta_{CP}\in\{90^\circ,270^\circ\}\) exactly.

\paragraph{Proof sketch.}
\(X=-\mathbf1\) forces a real mass matrix in the generation basis; the Fourier matrix of the triole turns
this into the \(\mu\)--\(\tau\) reflection, whose consequences are classical.  Checked on 2000 random
points of the four-dimensional solution space for each reflection type (script \texttt{k26\_test\_h8.py}).
\(\square\)

\subsection*{v4.8.0-r12 pre-registered tests}
\label{subsec:v480-r12-tests}

Data: NuFIT 6.0 (IC24 with SK atmospheric data, normal ordering unless stated), DESI DR2 with CMB for
\(\sum m_\nu<0.0642\) eV, KamLAND-Zen for \(m_{\beta\beta}\); Gaussian approximation with asymmetric
errors where a \(\chi^2\) is quoted.
\begin{center}
\small
\begin{tabular}{>{\raggedright\arraybackslash}p{0.06\textwidth}>{\raggedright\arraybackslash}p{0.10\textwidth}>{\raggedright\arraybackslash}p{0.36\textwidth}>{\raggedright\arraybackslash}p{0.38\textwidth}}
\toprule
Test & commit & hypothesis & result under the fixed criterion\\
\midrule
H2 & \texttt{434c5a8} & neutrinos are an \(r=1\) triole, signs free & compatible: normal ordering, \(\sum m_\nu=0.0592\pm0.0002\) eV; inverted ordering disfavoured (\(0.102\) eV); weakly discriminating\\
H3 & \texttt{e3b140e} & triole signs are Majorana CP parities & compatible, \(m_{\beta\beta}=2.7\pm0.7\) meV (\(\le3.5\) meV); not testable in the foreseeable future\\
H4 & \texttt{00d297e} & TM1 against TM2, equal prior & TM1 strongly preferred, \(\chi^2=4.4\) vs.\ \(14.7\) (2 d.o.f.); the \(\theta_{12}\) part was not blind\\
H5 & \texttt{16eff15} & polar angle of the preserved column is a Hopf angle of the Companion & failed: \(0/32\) canonical combinations compatible; free azimuth degenerate\\
H6 & \texttt{e1e2393} & list of algebraic selection criteria for the TM1 lift & main criterion (dihedral \(S_3\) with the triole) selects the \(\theta_{13}=0\) lift, i.e.\ speaks against TM1\\
H7 & \texttt{e6da7db} & Heisenberg compensator \(J\) as sign automorphism selects the lift & no selection for any \(J\); the candidates negating \(e_1\) and fixing \(e_7\) do not commute with the triole\\
H8 & \texttt{f6734ab} & Cayley--Dickson involution as CP symmetry & tension, \(\chi^2=8.8\) (3 d.o.f.); fixed predictions \(\sin^2\theta_{23}=\tfrac12\), \(\delta_{CP}=270^\circ\), with TM1 \(\sin^2\theta_{12}=0.318\)\\
\bottomrule
\end{tabular}
\end{center}

\paragraph{Reading.}
The octonionic structure already present in the Companion produces the known mixing patterns TM1, TM2
and \(\Delta(96)\) as candidates and, through the Cayley--Dickson involution, a definite CP structure.  It
does not select TM1: neither minimality, nor the cascade projections (they commute with all sign changes),
nor the compensator distinguishes the TM1 lift from the excluded one.  The predictions of H4 and H8 will be
tested by JUNO (\(\theta_{12}\)) and by DUNE and Hyper-Kamiokande (\(\theta_{23}\), \(\delta_{CP}\)).

\subsection*{v4.8.0-r12 reclassification of earlier PMNS statements}
\label{subsec:v480-r12-reclassification}

\paragraph{Proposition~\ref{prop:pmns-z3}.}
The value \(\arctan\sqrt2=54.74^\circ\) is not the maximal mixing angle; maximal mixing is \(45^\circ\).
The shifted value \(\arctan(\sqrt2)\,\delta_{\rm CL}=52.1^\circ\) gives \(\sin^2\theta_{23}=0.622\), against
\(0.470\) (normal ordering, \(\theta_{23}=43.3^\circ\)) about \(9\sigma\) in the Gaussian approximation and
about \(6\sigma\) against the inverted-ordering value \(0.550\).  The proposition and the
\(\theta_{23}^{\rm PMNS}\) line of the unified prediction table are reclassified as \emph{refuted}.  The
replacement statement of this layer is Theoremlet~r12.4 with the open selection problem.

\paragraph{Theorem~\ref{thm:op1op4-unified}.}
Its CKM and PMNS lines rest on the continuum-limit conjecture for \(\delta_{\rm CL}\), whose \(\beta\)-formula
was classified in layer r09 as a coincidence with an exponent chosen after the fact.  The CKM line inherits
that status; the PMNS line is refuted as above.  The CKM phase remark (\(\pi/3\) against \(\approx70^\circ\))
is unaffected and remains open.

\subsection*{v4.8.0-r12 ledger and guard update}
\label{subsec:v480-r12-ledger-guard}

\[
R_{\rm mix}^{v480\text{-}r12}=(0,0,1)
\]
for (protocols committed before computation; all seven hypotheses decided; selection principle for the
residual reflection still open).

The layer spends no guard:
\[
R_{\rm glob}^{(v480\text{-}r12,+)}=(0,1,1,1,1),
\qquad
\Gamma_{\rm guard}^{v480\text{-}r12}=(1,1,1,1,1),
\qquad
\beta_4=1/256\ \text{unspent}.
\]
It does not imply exact finite \(C5\) zero closure, exact all-readout terminality,
\(C6/T2\), a global photon-clock/RH theorem, or global all-readout closure.

% ============================================================================
\clearpage
\section*{v4.8.0-r13 Source-Only: JUNO 2026 Re-Check, Associator of the Generation Triple, and the Colour--Generation Assignment}
\addcontentsline{toc}{section}{v4.8.0-r13 Source-Only: JUNO 2026 Re-Check, Associator, Colour--Generation Assignment}
\label{sec:v480-r13-juno-associator-colour}
% ============================================================================

\paragraph{Purpose.}
This layer re-checks the fixed predictions of layer r12 against the JUNO data of September 2026, records two
further pre-registered items (H9, H10), states an obstruction for quark statements, and compares the local
line with the independent v4.9.0 research candidate.  Earlier layers are not rewritten.

\noindent\textbf{Source-only status.}
This v4.8.0-r13 layer is source-only.  It does not build a PDF, does not assign a new DOI,
and does not authorize a public release.  The v4.8.0 Version DOI
\href{https://doi.org/10.5281/zenodo.23365244}{10.5281/zenodo.23365244} and the Concept DOI
\href{https://doi.org/10.5281/zenodo.19210728}{10.5281/zenodo.19210728} are unchanged by it.

\subsection*{v4.8.0-r13 re-check against JUNO 2026}
\label{subsec:v480-r13-juno}

JUNO (arXiv:2609.37895, 207.2 live days) reports \(\sin^2\theta_{12}=0.3036\pm0.0064\),
\(\Delta m^2_{21}=(7.388\pm0.078)\times10^{-5}\) eV\(^2\), \(\Delta m^2_{31}=+2.509^{+0.027}_{-0.025}\times10^{-3}\)
eV\(^2\) for normal ordering, and prefers normal ordering by at least \(2.1\sigma\).  These data existed when H2--H8
were fixed but had not been used (the data basis was NuFIT~6.0); they are an independent re-check, not a test
fixed in advance of the measurement.  Data choice made before the computation: \(\sin^2\theta_{12}\) is the
inverse-variance mean of NuFIT~6.0 and JUNO, \(0.3046\pm0.0056\); for H2 the JUNO mass splittings; all other
inputs as in r12.  The decision criteria of r12 are unchanged.
\begin{center}
\small
\begin{tabular}{>{\raggedright\arraybackslash}p{0.10\textwidth}>{\raggedright\arraybackslash}p{0.25\textwidth}>{\raggedright\arraybackslash}p{0.25\textwidth}>{\raggedright\arraybackslash}p{0.28\textwidth}}
\toprule
Test & NuFIT 6.0 (r12) & with JUNO 2026 & verdict under the fixed criterion\\
\midrule
H2 & normal ordering, \(\sum m_\nu=0.0592\) eV & \(\sum m_\nu=0.0591\) eV; normal ordering preferred & compatible, ordering supported\\
H4 & \(\chi^2(\mathrm{TM1})=4.4\), \(\chi^2(\mathrm{TM2})=14.7\) & \(9.6\) and \(49.1\) & TM1 far ahead of TM2 but disfavoured on its own (\(\chi^2>9\), 2 d.o.f.)\\
H8 & \(\chi^2=8.8\), tension & \(\chi^2=14.0\) & disfavoured (\(\ge11.34\), 3 d.o.f.)\\
\bottomrule
\end{tabular}
\end{center}
The shift is driven by \(\theta_{12}\): TM1 requires \(\sin^2\theta_{12}=0.3182\), \(2.4\sigma\) above the combined
value.

\paragraph{H9 (note committed as \texttt{b75170d}).}
The relation \(|U_{e2}|^2=\beta_{\rm sym}^2\), found after the fact in the evaluation of H5, gives
\(\sin^2\theta_{12}=\beta_{\rm sym}^2/\cos^2\theta_{13}=0.30862\pm0.00018\) with
\(\beta_{\rm sym}=0.549348\) from \(\gamma_{\rm string}=\ln2/(\pi\sqrt3)\) and the inserted
\(\gamma_{\rm BI}=0.2375\).  Its present distance to JUNO (\(0.8\sigma\)) is disclosed and does not count.  The
test is the next JUNO result with \(\sigma\le0.004\): \(|z|<2\) compatible, \(|z|\ge3\) refuted.

\subsection*{v4.8.0-r13 associator of the generation triple}
\label{subsec:v480-r13-associator}

\paragraph{Theoremlet v4.8.0-r13.1.}
In the convention of layer r04, \([e_4,e_5,e_6]=(e_4e_5)e_6-e_4(e_5e_6)=-2e_7\).  The associator vanishes on Fano
lines and has norm \(2\) on non-collinear triples.  Hence every element of the stabilizer of \(+e_7\) fixes the
associator of the generation triple.

\paragraph{Proof sketch.}
From the multiplication of r04, \(e_4e_5=e_1\) and \((e_4e_5)e_6=e_1e_6=-e_7\), while \(e_5e_6=-e_3\) and
\(e_4(e_5e_6)=-e_4e_3=e_7\); the difference is \(-2e_7\).  All triples were also evaluated numerically
(script \texttt{k27\_test\_h10.py}).  An automorphism fixing \(e_7\)
maps \([e_4,e_5,e_6]\) to \([he_4,he_5,he_6]=h(-2e_7)=-2e_7\).  \(\square\)

\paragraph{H10 (note committed as \texttt{a6b5bcb}).}
Four associator criteria were fixed in advance for selecting the TM1 lift over the \(\theta_{13}=0\) lift.  None
separates the two, for any of the four lifts of the triole.  Together with r12 (minimality, cascade projections,
compensator) no internal selection principle for the TM1 lift has been found.

\subsection*{v4.8.0-r13 colour and generation}
\label{subsec:v480-r13-colour-generation}

Theorem~\ref{thm:three-generations} places the generations on \(e_4,e_5,e_6\), the triplet of
\(SU(3)\subset G_2\), and its proof gives each generation one colour index.  In the Dixon reading this \(SU(3)\)
is the colour group, so one octonionic triplet is made to carry both colour and generation.  Quarks require
\(3\times3\) such labels.  Consequences: (i) the octonionic flavour analysis of r12 is mathematically correct but
physically conditional on reading the triplet as generation space; read as colour, a broken flavour symmetry
inside \(SU(3)_c\) would break colour; (ii) quark values of the triole radius cannot be addressed before colour
and generation are separated.  In addition, the quark radii are scheme dependent at the \(2\)--\(4\%\) level
(\(r_u=1.290\) at \(M_Z\) against \(1.244\) at standard PDG scales; \(r_d=1.115\) against \(1.093\)), about
a hundred times the statistical error, so a structural match would carry no weight.  No fit is made.

\subsection*{v4.8.0-r13 comparison with the v4.9.0 research candidate}
\label{subsec:v480-r13-v49}

The repository's v4.9.0 research candidate (independent of this line) agrees with it on the Koide condition as an
equal-norm law without derivation, on rejecting the phase ansatz \(\pi-\sqrt2\arcsin\beta\), and on not
asserting a Riemann or GUE theorem.  It finds, independently, the positive phase chambers
\(|\theta-2\pi j/3|\le\pi/12\) with a parametrically light state near the boundary, the boundary of
Theoremlet~r12.1.  It admits only non-negative roots for the physical Koide ratio; H2 and H3 are therefore
reclassified as tests of the \emph{signed-root} Koide relation (for the H2 masses the non-negative-root ratio is
\(0.525\)).  The colour--generation assignment above concerns both lines.

\subsection*{v4.8.0-r13 ledger and guard update}
\label{subsec:v480-r13-ledger-guard}

\[
R_{\rm mix}^{v480\text{-}r13}=(0,0,1,1)
\]
for (protocols committed; hypotheses decided; selection principle open; colour--generation separation open).

The layer spends no guard:
\[
R_{\rm glob}^{(v480\text{-}r13,+)}=(0,1,1,1,1),
\qquad
\Gamma_{\rm guard}^{v480\text{-}r13}=(1,1,1,1,1),
\qquad
\beta_4=1/256\ \text{unspent}.
\]
It does not imply exact finite \(C5\) zero closure, exact all-readout terminality,
\(C6/T2\), a global photon-clock/RH theorem, or global all-readout closure.

% ============================================================================
\clearpage
\section*{v4.8.0-r14 Source-Only: Separating Colour and Generation}
\addcontentsline{toc}{section}{v4.8.0-r14 Source-Only: Separating Colour and Generation}
\label{sec:v480-r14-colour-generation}
% ============================================================================

\paragraph{Purpose.}
Layer r13 recorded that Theorem~\ref{thm:three-generations} puts colour and generation on the same octonionic
triplet.  This layer examines two ways of separating them, both built from structures already present in the
Companion, and records the consequences for layers r12 and r13.  It is model construction; no prediction is made
and no data are compared.  Earlier layers are not rewritten.

\noindent\textbf{Source-only status.}
This v4.8.0-r14 layer is source-only.  It does not build a PDF, does not assign a new DOI,
and does not authorize a public release.  The v4.8.0 Version DOI
\href{https://doi.org/10.5281/zenodo.23365244}{10.5281/zenodo.23365244} and the Concept DOI
\href{https://doi.org/10.5281/zenodo.19210728}{10.5281/zenodo.19210728} are unchanged by it.

\subsection*{v4.8.0-r14 requirements}
\label{subsec:v480-r14-requirements}

A separation should (a) keep colour as \(SU(3)\subset G_2\), the stabilizer of \(e_7\); (b) count generations by a
symmetry that commutes with \(G_2\), so that a broken generation symmetry does not break colour; and (c) provide
nine linearly independent coloured states, three colours times three generations.

\subsection*{v4.8.0-r14 route A: sedenions}
\label{subsec:v480-r14-sedenions}

\paragraph{Theoremlet v4.8.0-r14.1.}
Extend the Cayley--Dickson convention of layer r04 once more, \(\mathbb S=\OO\oplus\OO e_8\).  The maps
\(\psi\), a rotation by \(2\pi/3\) in each of the seven planes \((e_i,e_{i+8})\), and \(\varepsilon\colon
a+be_8\mapsto a-be_8\) are automorphisms of \(\mathbb S\), satisfy \(\psi^3=\varepsilon^2=1\) and
\(\varepsilon\psi\varepsilon=\psi^{-1}\), and commute with \(G_2\) acting diagonally.  On the fourteen plane
directions \(1+\psi+\psi^2=0\).  Under \(G_2\times S_3\),
\[
\mathbb S=(\mathbf1,\mathbf1)\oplus(\mathbf1,\mathrm{sgn})\oplus(\mathbf7,\mathbf2).
\]
Coloured states therefore carry only the two-dimensional representation of \(S_3\): the nine vectors
\(\psi^k(e_c)\), \(c=4,5,6\), \(k=0,1,2\), span a space of dimension \(6\).

\paragraph{Proof sketch.}
Automorphism, group relations and commutation were checked numerically for random arguments, for all \(1344\)
signed-permutation automorphisms and for random elements \(\exp D\) of the fourteen-dimensional derivation
algebra (script \texttt{k28\_sedenion\_generationen.py}); the colour stabilizer of \(e_7\) has dimension \(8\).
The traces \(\operatorname{tr}\psi=-5\), \(\operatorname{tr}\varepsilon=0\) fix the decomposition.  That
\(\operatorname{Aut}(\mathbb S)\) is exactly \(G_2\times S_3\) is taken from the literature (Gillard and
Gresnigt, arXiv:2306.13098, and references there), not re-proved.  \(\square\)

Route A separates colour from generation, but the triole relation \(z_1+z_2+z_3=0\) itself makes the three
copies of each coloured state linearly dependent.  It yields only two independent coloured generations.

\subsection*{v4.8.0-r14 route B: the triality form}
\label{subsec:v480-r14-triality}

\paragraph{Theoremlet v4.8.0-r14.2.}
On \(\OO\oplus\OO\oplus\OO\) the trilinear form \(t(x,y,z)=\operatorname{Re}\bigl((xy)z\bigr)\) is invariant under the
cyclic shift \((x,y,z)\mapsto(y,z,x)\), under \((x,y,z)\mapsto(\bar x,\bar z,\bar y)\), which together generate
\(S_3\), and under \(G_2\) acting diagonally, which commutes with this \(S_3\).  The three copies of the colour
triplet \((e_4,e_5,e_6)\) give nine linearly independent coloured states.

\paragraph{Proof sketch.}
The real part of the associator vanishes and \(\operatorname{Re}(ab)=\operatorname{Re}(ba)\), so
\(\operatorname{Re}((xy)z)=\operatorname{Re}(x(yz))=\operatorname{Re}((yz)x)\).  For the reflection,
\(\operatorname{Re}((\bar x\bar z)\bar y)=\operatorname{Re}(y(zx))=\operatorname{Re}((xy)z)\).  Elements of
\(G_2\) preserve multiplication and conjugation.  Independence holds because the three copies lie in different
summands; all statements were also checked numerically.  \(\square\)

\paragraph{Proposed assignment.}
Generation index = the three triality slots; colour index = the \(SU(3)\) triplet \(e_4,e_5,e_6\) within each
slot.  This matches the Companion's identification of \(S_2\cong S^8=\OO P^1\) with three generations, and the
fact that \(G_2\) is the subgroup of \(\mathrm{Spin}(8)\) fixed by triality.  The assignment is proposed, not
derived; a native carrier for the three slots is still to be supplied.

\subsection*{v4.8.0-r14 consequences}
\label{subsec:v480-r14-consequences}

\begin{itemize}
\item After the separation, the order-\(96\) stabilizer of Theoremlet~r12.2 lies inside the colour group and
cannot act as a flavour symmetry.  Theoremlets r12.2--r12.4 remain correct mathematics but lose their reading
as lepton-mixing statements; H4 and H8 were already disfavoured by the JUNO 2026 re-check of layer r13.
\item A pure generation \(S_3\) with the triole \(\mathbb Z_3\) for charged leptons and a transposition for
neutrinos preserves the column \((0,\tfrac12,\tfrac12)\), i.e.\ \(\theta_{13}=0\), which is excluded.  This is
the minimality outcome of H6.  Realistic mixing requires further breaking, from which no prediction follows yet.
\item The triole form \(t(1,1,1)+\sqrt2\,\mathrm{Re}(c(1,\omega,\omega^2))\) of layer r10 is the decomposition of
generation space under \(S_3\) into the symmetric direction and the doublet; \(K=2/3\) remains the equal-norm
condition between them and is still not derived.
\end{itemize}

\subsection*{v4.8.0-r14 ledger and guard update}
\label{subsec:v480-r14-ledger-guard}

\[
R_{\rm cg}^{v480\text{-}r14}=(0,1,1)
\]
for (separation consistent with \(G_2\) achieved by route B; native carrier of the triality slots open; mixing
mechanism open).

The layer spends no guard:
\[
R_{\rm glob}^{(v480\text{-}r14,+)}=(0,1,1,1,1),
\qquad
\Gamma_{\rm guard}^{v480\text{-}r14}=(1,1,1,1,1),
\qquad
\beta_4=1/256\ \text{unspent}.
\]
It does not imply exact finite \(C5\) zero closure, exact all-readout terminality,
\(C6/T2\), a global photon-clock/RH theorem, or global all-readout closure.

% ============================================================================
\clearpage
\section*{v4.8.0-r15 Source-Only: The Octonionic Hopf Fibration as Carrier of the Triality Slots}
\addcontentsline{toc}{section}{v4.8.0-r15 Source-Only: Octonionic Hopf Fibration as Carrier of the Triality Slots}
\label{sec:v480-r15-hopf-triality-carrier}
% ============================================================================

\paragraph{Purpose.}
Layer r14 proposed the three triality slots of \(\operatorname{Re}((xy)z)\) as generation index but left their
carrier open.  This layer shows that the octonionic Hopf fibration of Theorem~\ref{thm:hopf}, which the Companion
already attaches to \(S_2\cong\OO P^1\) and to three generations, carries the slots.  Model construction only; no
prediction is made and no data are compared.  Earlier layers are not rewritten.

\noindent\textbf{Source-only status.}
This v4.8.0-r15 layer is source-only.  It does not build a PDF, does not assign a new DOI,
and does not authorize a public release.  The v4.8.0 Version DOI
\href{https://doi.org/10.5281/zenodo.23365244}{10.5281/zenodo.23365244} and the Concept DOI
\href{https://doi.org/10.5281/zenodo.19210728}{10.5281/zenodo.19210728} are unchanged by it.

\subsection*{v4.8.0-r15 unit triples and the Hopf map}
\label{subsec:v480-r15-triples}

\paragraph{Theoremlet v4.8.0-r15.1.}
Let \(\mathcal T=\{(a,b,c)\in(S^7)^3:(ab)c=1\}\).  Then \(\mathcal T\) is a fourteen-dimensional manifold
(\(a,b\) free, \(c=\overline{ab}\)), closed under the cyclic shift \((a,b,c)\mapsto(b,c,a)\) and under
\((a,b,c)\mapsto(\bar a,\bar c,\bar b)\), and invariant under \(G_2\) acting diagonally.  For the Hopf map
\(h(x,y)=(|x|^2-|y|^2,\,2x\bar y)\colon S^{15}\to S^8\), the points over the equator \(|x|=|y|=1/\sqrt2\) give
unit triples \((\sqrt2\,x,\sqrt2\,\bar y,\bar v)\in\mathcal T\) with \(v=2x\bar y\), and the fibre over \(v\) is
\(\{(a,\bar a v)\}\cong S^7\).

\paragraph{Proof sketch.}
\(c=\overline{ab}\) gives \((ab)c=|ab|^2=1\).  By alternativity \(b(\bar b\bar a)=\bar a\), so \(bc=\bar a\) and
\((bc)a=1\); the remaining cyclic and reflected identities follow in the same way.  Elements of \(G_2\) preserve
products and conjugation.  The Hopf statements are the norm identity \((|x|^2-|y|^2)^2+4|x|^2|y|^2=1\) and
\(a(\bar a v)=v\).  All identities were also checked numerically (script \texttt{k29\_hopf\_triality\_traeger.py}).
\(\square\)

\paragraph{Theoremlet v4.8.0-r15.2 (symmetric point).}
\((a,a,a)\in\mathcal T\) if and only if \(a^3=1\), i.e.\ \(a=\exp(2\pi u/3)\) for an imaginary unit \(u\).  For
\(u=e_7\) the point is \(\omega\) in \(\mathrm{span}(1,e_7)\), the complex plane at the end of the cascade, and
its stabilizer in \(G_2\) is the colour group \(SU(3)\) (dimension \(8\)).

\paragraph{Proof sketch.}
\(a\) and its powers lie in the commutative subalgebra \(\mathrm{span}(1,u)\cong\mathbb C\); the stabilizer of
\(a\) in \(G_2\) is the stabilizer of \(u\).  The dimension was computed from the derivation algebra.  \(\square\)

\subsection*{v4.8.0-r15 reading and a check}
\label{subsec:v480-r15-reading}

The two components of the Hopf total space and the base point are the three triality slots; the triole is their
cyclic symmetry, and \(G_2\) acts as the colour side.  Choosing the symmetric generation point fixes \(u\) and breaks
\(G_2\) to the colour \(SU(3)\).  This gives a natural carrier for the assignment of r14; it does not derive the
existence of three generations.

Check: in \(\mathrm{span}(1,e_7)\) the condition \((ab)c=1\) becomes \(\alpha_1+\alpha_2+\alpha_3\equiv0\pmod{2\pi}\)
for the phases.  Reading the Koide phases \(\delta+2\pi k/3\) of layer r10 as these phases requires
\(3\delta\equiv0\), against \(\delta=0.2222\) for the charged leptons.  The direct reading fails.  That
\(3\delta\approx2/3\) is the known value \(2/9\cdot3\), whose generality was refuted in r11; it is not counted.  How
a mass hierarchy and \(K=2/3\) could arise from deviations from the symmetric triple remains open.

\subsection*{v4.8.0-r15 ledger and guard update}
\label{subsec:v480-r15-ledger-guard}

\[
R_{\rm cg}^{v480\text{-}r15}=(0,0,1)
\]
for (separation consistent with \(G_2\); carrier of the triality slots identified; mass mechanism open).

The layer spends no guard:
\[
R_{\rm glob}^{(v480\text{-}r15,+)}=(0,1,1,1,1),
\qquad
\Gamma_{\rm guard}^{v480\text{-}r15}=(1,1,1,1,1),
\qquad
\beta_4=1/256\ \text{unspent}.
\]
It does not imply exact finite \(C5\) zero closure, exact all-readout terminality,
\(C6/T2\), a global photon-clock/RH theorem, or global all-readout closure.

% ============================================================================
\clearpage
\section*{v4.8.0-r16 Source-Only: Pre-Registered Test of Masses from Invariants on the Triple Space}
\addcontentsline{toc}{section}{v4.8.0-r16 Source-Only: Pre-Registered Test of Masses from Invariants on the Triple Space}
\label{sec:v480-r16-triple-space-masses}
% ============================================================================

\paragraph{Purpose.}
Layer r15 left open whether a mass hierarchy can arise on the space of triality triples.  This layer records the
pre-registered test H11 (note \texttt{Vorab\_Notiz\_H11\_Massen\_Tripelraum.md}, commit \texttt{8e5eaa9}, fixed
before computing): are the masses the slot norms at critical points of a natural invariant function?  Earlier
layers are not rewritten.

\noindent\textbf{Source-only status.}
This v4.8.0-r16 layer is source-only.  It does not build a PDF, does not assign a new DOI,
and does not authorize a public release.  The v4.8.0 Version DOI
\href{https://doi.org/10.5281/zenodo.23365244}{10.5281/zenodo.23365244} and the Concept DOI
\href{https://doi.org/10.5281/zenodo.19210728}{10.5281/zenodo.19210728} are unchanged by it.

\subsection*{v4.8.0-r16 setting and hypotheses}
\label{subsec:v480-r16-setting}

For \((a,b,c)\in\OO^3\) put \(n_k=|{\rm slot}\ k|\), \(S=\sum n_k^2\), \(Q=\sum n_k^4\), \(P=n_1n_2n_3\) and
\(t=\operatorname{Re}((ab)c)\), \(|t|\le P\).  These generate the invariants of \(G_2\) (diagonal)
\(\times\,S_3\) up to degree four.  The fixed reading is \(\sqrt{m_k}\propto n_k\), so
\(K=\sum n_k^2/(\sum n_k)^2\).
\begin{itemize}
\item[H11a:] the masses are the slot norms at critical points of \(t\) on \(\{S=3\}\);
\item[H11b:] the masses are the slot norms at minima of
\(V=\alpha S+\beta S^2+\gamma Q-\kappa t\), bounded below.
\end{itemize}
Criterion (fixed): refuted if no critical point, respectively no minimum, has three distinct non-zero norms.

\subsection*{v4.8.0-r16 results}
\label{subsec:v480-r16-results}

\paragraph{Theoremlet v4.8.0-r16.1.}
At every critical point of \(V\) with all \(n_k\neq0\) and \(t=P\), each \(x_k=n_k^2\) solves the same equation
\[
(2\alpha+4\beta S)\,x+4\gamma\,x^2=\kappa P .
\]
Hence the slot norms take at most two distinct values.  The same holds for the critical points of \(t\) on
\(\{S=3\}\), which have equal norms (\(K=1/3\)) or two empty slots (\(K=1\)).

\paragraph{Proof sketch.}
\(\partial V/\partial n_k=0\) multiplied by \(n_k\) gives the displayed equation, since \(n_k\,\partial P/\partial
n_k=P\) for every \(k\).  A quadratic has at most two roots.  For H11a the Lagrange conditions give
\(|bc|=\lambda|a|\) and its cyclic images.  \(\square\)

Numerically (script \texttt{k30\_test\_h11.py}) the global minima on a grid of 210 bounded parameter points are:
180 with equal norms, 25 with one occupied slot, 5 flat (\(\gamma=\kappa=0\), a whole sphere of minima), none with
three distinct norms.  One deviation from the note is disclosed: the relative tolerance for ``distinct'' was raised
from \(10^{-6}\) to \(10^{-4}\) because the smaller value counted numerical noise in favour of the hypothesis.  The
verdict does not depend on it, since Theoremlet r16.1 excludes three distinct norms analytically.

\paragraph{Verdict.}
H11a and H11b are refuted under the fixed criterion.  Three distinct norms need a term of degree at least six,
which makes the equation cubic; this is necessary, not shown to be sufficient (no such minimum appeared in a
tested sextic grid).  Any such extension adds at least two parameters, so the two charged-lepton ratios could be
fitted but not predicted, and \(K=2/3\) would be a condition on the parameters rather than a consequence.  The
triole gives the form of the Koide relation (r10), not its dynamics.

\subsection*{v4.8.0-r16 ledger and guard update}
\label{subsec:v480-r16-ledger-guard}

\[
R_{\rm cg}^{v480\text{-}r16}=(0,0,1)
\]
for (separation consistent with \(G_2\); carrier of the triality slots identified; mass mechanism open, with
invariants up to degree four excluded).

The layer spends no guard:
\[
R_{\rm glob}^{(v480\text{-}r16,+)}=(0,1,1,1,1),
\qquad
\Gamma_{\rm guard}^{v480\text{-}r16}=(1,1,1,1,1),
\qquad
\beta_4=1/256\ \text{unspent}.
\]
It does not imply exact finite \(C5\) zero closure, exact all-readout terminality,
\(C6/T2\), a global photon-clock/RH theorem, or global all-readout closure.

% ============================================================================
\clearpage
\section*{v4.8.0-r17 Source-Only: Time Direction of the Cosmological Fractions}
\addcontentsline{toc}{section}{v4.8.0-r17 Source-Only: Time Direction of the Cosmological Fractions}
\label{sec:v480-r17-cosmological-time-direction}
% ============================================================================

\paragraph{Purpose.}
Layer r03 examined the assignment \(\Omega_\Lambda=1-\bsym^2\) of Theorem~\ref{thm:dark-energy} at the present
epoch (\S\ref{subsec:v480-r03-p4-recheck}) and the provenance of \(\bsym\)
(\S\ref{subsec:v480-r03-frame-parameter-provenance}).  Two hypothesis-grade statements that make the same
assignment time-dependent, Proposition~\ref{prop:H0-tension} and Proposition~\ref{prop:saturation}, have not been
checked by any layer since r85 and do not appear in the registers of \S\ref{sec:v480-r09-p2-p3-reclassification}.
This layer compares them with the measured expansion history.  It is a consistency check of existing statements,
not a new prediction, and therefore carries no pre-registration note.  The computation is
\texttt{llm/k31\_kosmologie\_vorzeichen.py}.  Earlier layers are not rewritten.

\noindent\textbf{Source-only status.}
This v4.8.0-r17 layer is source-only.  It does not build a PDF, does not assign a new DOI,
and does not authorize a public release.  The v4.8.0 Version DOI
\href{https://doi.org/10.5281/zenodo.23365244}{10.5281/zenodo.23365244} and the Concept DOI
\href{https://doi.org/10.5281/zenodo.19210728}{10.5281/zenodo.19210728} are unchanged by it.

\subsection*{v4.8.0-r17 the three statements}
\label{subsec:v480-r17-statements}

\begin{enumerate}[leftmargin=2em]
\item Theorem~\ref{thm:dark-energy}: \(\Omega_\Lambda=P(S^8\to S^2)=1-\bsym^2\) and
\(\Omega_M=P(S^8\to S^4)=\bsym^2\).
\item Proposition~\ref{prop:H0-tension}: \(\bsym\) evolves, with \(\bsym(z=1100)\approx0.527\) and
\(\bsym(z=0)\approx0.5493\); ``the universe moves continuously away from the photon frame''.
\item Proposition~\ref{prop:saturation}, stated there as ``the consistent continuation'' of item~2:
\(\bsym(t)\nearrow1\), \(P(S^8\to S^4)=\bsym^2\to1\), ``the universe becomes fully matter dominated'', and the
expansion stops.
\end{enumerate}
Items~2 and~3 use item~1 at all epochs, not only today; item~3 does so explicitly through \(\bsym^2\to1\).

\subsection*{v4.8.0-r17 comparison with the expansion history}
\label{subsec:v480-r17-comparison}

In flat \(\Lambda\)CDM with the Planck 2018 values \(\Omega_M=0.3111\), \(\Omega_r=9.1\times10^{-5}\), the
\(\Lambda\) fraction at redshift \(z\) is
\(\Omega_\Lambda(z)=\Omega_\Lambda/\bigl(\Omega_r(1+z)^4+\Omega_M(1+z)^3+\Omega_\Lambda\bigr)\).  Under item~1 the
value of \(\bsym\) that reproduces it is \(\sqrt{1-\Omega_\Lambda(z)}\):
\begin{center}
\small
\begin{tabular}{@{}rrr@{}}
\toprule
\(z\) & \(\Omega_\Lambda(z)\) & required \(\bsym\)\\
\midrule
\(-0.5\) & \(0.947\) & \(0.231\)\\
\(0\) & \(0.689\) & \(0.558\)\\
\(0.5\) & \(0.396\) & \(0.777\)\\
\(1\) & \(0.217\) & \(0.885\)\\
\(2\) & \(0.0757\) & \(0.961\)\\
\(5\) & \(0.0101\) & \(0.995\)\\
\(1100\) & \(1.3\times10^{-9}\) & \(1.000\)\\
\bottomrule
\end{tabular}
\end{center}
Proposition~\ref{prop:H0-tension} gives instead \(\Omega_\Lambda(1100)=1-0.527^2=0.722\), larger than the
\(\Lambda\)CDM value by a factor \(5.8\times10^{8}\).

\paragraph{Theoremlet v4.8.0-r17.1.}
Under the assignment of Theorem~\ref{thm:dark-energy} held at all epochs, the measured expansion history requires
\(\mathrm d\bsym/\mathrm dt<0\): \(\bsym\) decreases from \(\bsym\approx1\) at recombination to \(0.558\) today and
tends to \(0\) in the \(\Lambda\)-dominated future.  Propositions~\ref{prop:H0-tension} and~\ref{prop:saturation}
assume \(\mathrm d\bsym/\mathrm dt>0\) and an end state \(\bsym\to1\).  The three statements together describe a
universe that starts \(\Lambda\)-dominated and ends matter-dominated, the reverse of the observed history.  They are
therefore not jointly tenable.

\paragraph{Proof sketch.}
\(\Omega_\Lambda(z)\) is monotonically decreasing in \(z\) for \(\Omega_M,\Omega_\Lambda>0\), so
\(\sqrt{1-\Omega_\Lambda(z)}\) is increasing in \(z\), i.e.\ decreasing in cosmic time.  The table evaluates it.
Items~2 and~3 state the opposite monotonicity.  \(\square\)

\paragraph{What the comparison does not decide.}
The present-epoch agreement \(1-\bsym^2=0.698\) against \(\Omega_\Lambda=0.689\) (\(1.4\%\)) is untouched, as are
the caveats of r03 F3--F5 and F11.  The check uses \(\Lambda\)CDM as the description of the measured history; a
different background would change the numbers, not the sign, because the ratio of matter to dark-energy density
scales as \(a^{3w}\) and falls with expansion for every dark-energy equation of state \(w<0\).  This also removes a possible affinity with the
current evidence for evolving dark energy: that evidence concerns the dark-energy density, while
Proposition~\ref{prop:H0-tension} concerns the fraction, which increases with time in every such model.

\paragraph{Ways out, recorded and not chosen.}
\begin{enumerate}[leftmargin=2em]
\item[(a)] Theorem~\ref{thm:dark-energy} holds only at the present epoch.  Then
Propositions~\ref{prop:H0-tension} and~\ref{prop:saturation} lose their basis, and the \(1.4\%\) agreement becomes
a statement about one epoch, which r03 F11 already qualifies.
\item[(b)] \(\bsym\) runs the other way, from \(\approx1\) at recombination downward.  Then the early universe sits at
what Proposition~\ref{prop:saturation} calls the tachyon limit, and Proposition~\ref{prop:H0-tension} would need the
opposite sign of \(\Delta\bsym\), which it would have to re-derive.
\item[(c)] The cascade probabilities are not the cosmological fractions at other epochs, but some other function of
them.  Then item~1 needs a time-dependent bridge that the source does not supply.
\end{enumerate}
Which way is taken is left to the author.

\subsection*{v4.8.0-r17 register entries}
\label{subsec:v480-r17-register}

\begin{center}
\small
\begin{tabular}{@{}>{\raggedright\arraybackslash}p{0.34\textwidth}>{\raggedright\arraybackslash}p{0.58\textwidth}@{}}
\toprule
statement & status after r17\\
\midrule
Thm.~\ref{thm:dark-energy} & present-epoch value \(1.4\%\) from Planck; not tenable at all epochs together with Props.~\ref{prop:H0-tension}, \ref{prop:saturation} (r17.1)\\
Prop.~\ref{prop:H0-tension} & sign of \(\mathrm d\bsym/\mathrm dt\) opposite to the one required by Thm.~\ref{thm:dark-energy}; \(\Omega_\Lambda(1100)\) off by \(5.8\times10^8\) (r17)\\
Prop.~\ref{prop:saturation} & end state \(\Omega_M\to1\) reverses the observed trend \(\Omega_\Lambda\to1\) (r17)\\
\bottomrule
\end{tabular}
\end{center}
These entries extend the register of \S\ref{subsec:v480-r09-results}.

\subsection*{v4.8.0-r17 ledger and guard update}
\label{subsec:v480-r17-ledger-guard}

\[
R_{\rm cos}^{v480\text{-}r17}=(0,0,1)
\]
for (present-epoch agreement recorded, time direction checked against the expansion history, choice among
(a)--(c) open).

The layer spends no guard:
\[
R_{\rm glob}^{(v480\text{-}r17,+)}=(0,1,1,1,1),
\qquad
\Gamma_{\rm guard}^{v480\text{-}r17}=(1,1,1,1,1),
\qquad
\beta_4=1/256\ \text{unspent}.
\]
It does not imply exact finite \(C5\) zero closure, exact all-readout terminality,
\(C6/T2\), a global photon-clock/RH theorem, or global all-readout closure.

% ============================================================================
\clearpage
\section*{v4.8.0-r18 Source-Only: Three-Frame Typing of the Frame Offset and Correction of r17}
\addcontentsline{toc}{section}{v4.8.0-r18 Source-Only: Three-Frame Typing of the Frame Offset and Correction of r17}
\label{sec:v480-r18-three-frame-typing}
% ============================================================================

\paragraph{Purpose.}
The single symbol \(\bsym\) is used in 261 lines of the main body in at least three roles: as a frame factor in
the conversion of constants, as a branching probability of the cascade, and as a Hopf angle
(\texttt{llm/k32\_bsym\_rollen.py}).  This layer types the frame offset over the three frames of the triolic
system, records three findings that the single symbol hides, narrows the register entries of r17, and records a
reading of why the offsets coincide today.  It contains arithmetic and typing, no new prediction.  Earlier layers
are not rewritten.

\noindent\textbf{Source-only status.}
This v4.8.0-r18 layer is source-only.  It does not build a PDF, does not assign a new DOI,
and does not authorize a public release.  The v4.8.0 Version DOI
\href{https://doi.org/10.5281/zenodo.23365244}{10.5281/zenodo.23365244} and the Concept DOI
\href{https://doi.org/10.5281/zenodo.19210728}{10.5281/zenodo.19210728} are unchanged by it.

\subsection*{v4.8.0-r18 correction of r17}
\label{subsec:v480-r18-r17-correction}

Proposition~\ref{prop:H0-tension} does not mention the cosmological fractions.  It states that \(\bsym\) drifts
and that the \(H_0\) discrepancy is a frame artefact; the drift acts through the frame factor.  (The ratio
\(\gL(z{=}0)/\gL(z{=}1100)=1.8642/1.7690=1.054\) raised to the power \(3/2\) gives \(1.082\) against the quoted
\(1.083\); the exponent is a reconstruction, the source does not state it.)  The value
\(\Omega_\Lambda(1100)=0.722\) of \S\ref{subsec:v480-r17-comparison} arises only when the drifting offset is inserted
into Theorem~\ref{thm:dark-energy}.  The one statement of the source that makes this insertion is item~(2) of
Proposition~\ref{prop:saturation}, \(P(S^8\to S^4)=\bsym^2\to1\).

Theoremlet r17.1 stands, because it is conditional on the assignment holding at all epochs.  The register entries of
\S\ref{subsec:v480-r17-register} are narrowed as follows: Proposition~\ref{prop:H0-tension} is not refuted by r17;
the conflict with the expansion history sits in item~(2) of Proposition~\ref{prop:saturation} and in the tacit
identification of two offsets that \S\ref{subsec:v480-r18-reading} separates.  Item~(1) of
Proposition~\ref{prop:saturation} uses the frame factor only and is not affected.  The word ``permanent'' in
Theorem~\ref{thm:dark-offset} means non-vanishing (``we can never fully reach \(\zminus=0\)''); it does not
conflict with a slow drift.

\subsection*{v4.8.0-r18 three frames and their cross terms}
\label{subsec:v480-r18-cross-terms}

Write \(\eta=\operatorname{artanh}\bsym=0.6174\) for the rapidity of the offset.  By
Theorem~\ref{thm:symmetric-frame} matter (\(+\)), photon (\(0\)) and tachyon (\(-\)) sit at rapidities
\(\eta_+=\eta\), \(\eta_0=0\), \(\eta_-=-\eta\).  Rapidities add, so the offset of frame \(a\) seen from frame \(b\) is
\(\eta_{ab}=\eta_a-\eta_b\), with velocity \(\tanh\eta_{ab}\) and Lorentz factor \(\cosh\eta_{ab}\):
\begin{center}
\small
\begin{tabular}{@{}c|ccc@{}}
\toprule
from\,\(\backslash\)\,to & \(+\) & \(0\) & \(-\)\\
\midrule
\(+\) & \(0\) & \(0.5493,\ \gamma=1.197\) & \(0.8440,\ \gamma=1.864\)\\
\(0\) & \(-0.5493,\ \gamma=1.197\) & \(0\) & \(0.5493,\ \gamma=1.197\)\\
\(-\) & \(-0.8440,\ \gamma=1.864\) & \(-0.5493,\ \gamma=1.197\) & \(0\)\\
\bottomrule
\end{tabular}
\end{center}

\paragraph{Theoremlet v4.8.0-r18.1.}
Under the symmetry \(\eta_-=-\eta_+\) of Theorem~\ref{thm:symmetric-frame} the nine cross terms take exactly two
non-trivial magnitudes: the neighbour term (\(+0\), \(0-\)) with \(\cosh\eta=1.197\) and the opposite term (\(+-\))
with \(\cosh2\eta=(1+\bsym^2)/(1-\bsym^2)=\gL=1.864\).  The frame factor \(\gL\) is the opposite term.

\paragraph{Proof sketch.}
\(\cosh2\eta=(1+\tanh^2\eta)/(1-\tanh^2\eta)\) with \(\tanh\eta=\bsym\); this is the formula of
Theorem~\ref{thm:symmetric-frame}, which names \(\gL\) the Lorentz factor between the matter and tachyon frames.
\(\square\)

\subsection*{v4.8.0-r18 findings}
\label{subsec:v480-r18-findings}

\paragraph{(a) The conversion of constants uses the opposite term for a neighbour pair.}
The definition of the frame-insertion count in \S\ref{sec:phase-constants} converts matter to photon by
\(X^{(\mathrm{photon})}=\gL^{-k(X)}X^{(\mathrm{matter})}\), and the proof of Theorem~\ref{thm:frame-offset} reads
\(\gL\) as a Lorentz factor between the matter frame and the photon frame, with \(\beta=0.844\).  By r18.1 the
matter--photon pair is a neighbour pair, with factor \(1.197\); \(1.864\) belongs to matter--tachyon.  Since
\(\cosh2\eta\neq(\cosh\eta)^n\) for integer \(n\) (\((\cosh\eta)^2=1.432\)), the difference cannot be absorbed into the
counts \(k\).  Either the conversion is meant to pass through the tachyon frame, which would have to be stated, or
every frame-corrected constant of \S\ref{sec:phase-constants} uses the wrong cross term.  r18 does not choose.

\paragraph{(b) The observer position \(\zminus=\beta/2=0.422\).}
Theorem~\ref{thm:frame-offset} places the observer at \(\zminus=\beta/2\approx0.422\), halving the opposite offset
\(0.844\) in velocity.  Halving in rapidity gives \(\tanh\eta=0.5493=\bsym\), the position of
Theorem~\ref{thm:symmetric-frame}.  The value \(0.422\) is a halving in the wrong variable; the two theorems agree
once the halving is done in rapidity.

\paragraph{(c) The cascade split is the neighbour term.}
Theorem~\ref{thm:dark-energy} has \(\Omega_M=\tanh^2\eta=0.3018\) and \(\Omega_\Lambda=\operatorname{sech}^2\eta=0.6982\),
which sum to one identically.  The split is therefore internally consistent.  Open: the cascade starts in the
tachyon stratum \(S^8\), yet uses the neighbour term \(+0\) and no term of the \(-\) row.

\subsection*{v4.8.0-r18 a reading of the coincidence, offered and not adopted}
\label{subsec:v480-r18-reading}

A single \(\bsym\) is the statement \(\eta_-=-\eta_+\).  If the matter frame drifts, as
Proposition~\ref{prop:H0-tension} has it, while the tachyon frame stays at a structural value \(-\eta_s\), the symmetry
is broken and the cross terms become three distinct numbers, \(\eta_+(t)\), \(\eta_s\) and \(\eta_+(t)+\eta_s\).
A structural offset (cascade side) and an observed offset (frame side) then coincide exactly at the symmetric point
\(\eta_+=\eta_s\).

The source supplies a reason for that point.  In Theorem~\ref{thm:spring} the net force on the photon,
\(k(z^-_T-z^-_\gamma)+k(z^-_M-z^-_\gamma)\), vanishes for \(z^-_\gamma=0\) only if \(z^-_T=-z^-_M\): the symmetric
configuration is the mechanical equilibrium of the triple.  Under this reading ``the offsets agree today'' means
``the triple is near its equilibrium today'', and the gap between \(0.5493\) (from \(\gL\)) and \(0.5578\) (from
\(\Omega_M\), Remark~\ref{rem:omega-bsym}), \(1.5\%\), would be a departure from symmetry rather than the lattice
correction \(\delta_{\mathrm{CL}}\) to which the remark attributes it.  The system of Theorem~\ref{thm:spring} is
undamped: it supplies an equilibrium, not a relaxation towards it.

Which measured value belongs to which cross term under broken symmetry is not decided here.  Any such assignment is a
prediction and requires a pre-registration note before it is computed.

\subsection*{v4.8.0-r18 register entries}
\label{subsec:v480-r18-register}

\begin{center}
\small
\begin{tabular}{@{}>{\raggedright\arraybackslash}p{0.34\textwidth}>{\raggedright\arraybackslash}p{0.58\textwidth}@{}}
\toprule
statement & status after r18\\
\midrule
Prop.~\ref{prop:H0-tension} & not refuted by r17; uses the offset only through the frame factor (r18)\\
Prop.~\ref{prop:saturation} & item (2) carries the conflict of r17.1; item (1) unaffected (r18)\\
\(\gL\) in \S\ref{sec:phase-constants}, Thm.~\ref{thm:frame-offset} & opposite cross term (\(+-\)) used for the neighbour pair matter--photon (r18 (a))\\
Thm.~\ref{thm:frame-offset}, \(\zminus=0.422\) & halving in velocity; in rapidity \(\zminus=\bsym\) (r18 (b))\\
Thm.~\ref{thm:dark-energy} & split \(\tanh^2\eta+\operatorname{sech}^2\eta=1\) consistent; choice of the \(+0\) term for a cascade from \(S^8\) open (r18 (c))\\
\bottomrule
\end{tabular}
\end{center}

\subsection*{v4.8.0-r18 ledger and guard update}
\label{subsec:v480-r18-ledger-guard}

\[
R_{\rm fr}^{v480\text{-}r18}=(0,0,1)
\]
for (three-frame cross terms typed, findings (a)--(c) recorded, assignment of the measured values to cross terms
open).

The layer spends no guard:
\[
R_{\rm glob}^{(v480\text{-}r18,+)}=(0,1,1,1,1),
\qquad
\Gamma_{\rm guard}^{v480\text{-}r18}=(1,1,1,1,1),
\qquad
\beta_4=1/256\ \text{unspent}.
\]
It does not imply exact finite \(C5\) zero closure, exact all-readout terminality,
\(C6/T2\), a global photon-clock/RH theorem, or global all-readout closure.

% ============================================================================
\clearpage
\section*{v4.8.0-r19 Source-Only: Frame Cubes, Transition Group and Holonomy}
\addcontentsline{toc}{section}{v4.8.0-r19 Source-Only: Frame Cubes, Transition Group and Holonomy}
\label{sec:v480-r19-frame-cubes-holonomy}
% ============================================================================

\paragraph{Purpose.}
Layer r18 typed the frame offset by cross terms between the three frames.  Two working results followed outside the
source: the pre-registered test H12 (\texttt{llm/Ergebnis\_H12\_Drei\_Frame\_Zuordnung.md}) and a backward analysis
with a readout index (\texttt{llm/k34\_rueckwaerts\_3x3x3.py}, map entry K39).  This layer records both, asks what a
family of readout cubes linked by transition maps adds, and compares the transition group of the earlier reading with
that of Reading~B (\S\ref{subsec:v480-r06-readings}).  It contains typing and arithmetic, no new prediction.  Earlier
layers are not rewritten.

\noindent\textbf{Source-only status.}
This v4.8.0-r19 layer is source-only.  It does not build a PDF, does not assign a new DOI,
and does not authorize a public release.  The v4.8.0 Version DOI
\href{https://doi.org/10.5281/zenodo.23365244}{10.5281/zenodo.23365244} and the Concept DOI
\href{https://doi.org/10.5281/zenodo.19210728}{10.5281/zenodo.19210728} are unchanged by it.

\subsection*{v4.8.0-r19 records of H12 and of the backward analysis}
\label{subsec:v480-r19-records}

\paragraph{H12.}
With \(\cosh(\eta_++\eta_s)=\gL\), \(\Omega_M=\tanh^2\eta_s\) (Planck 2018 with BAO) and
\(\sin^2\theta_{12}=\tanh^2\eta_+/\cos^2\theta_{13}\), the assignment gives \(\beta_+=0.5408\), \(\beta_s=0.5578\) and
\(\sin^2\theta_{12}=0.2991\pm0.0057\), against \(0.3086\) for the symmetric relation H9.  At the precision fixed for
the next JUNO release the two are \(1.36\sigma\) apart; by the fixed criterion they cannot be separated, and the
result is recorded as no test.  The limit is set by the uncertainty of \(\Omega_M\), not by JUNO.

\paragraph{Backward analysis (does not count as a test).}
Let \(T_{abc}\) be the offset of frame \(a\) against frame \(b\), read from frame \(c\), starting from the Hopf value
\(\bsym=1/\sqrt3\).  Pure kinematics makes \(T_{abc}=\eta_a-\eta_b\) independent of \(c\), since a boost into \(c\)
shifts all rapidities equally.  Of the readout terms tried, those depending on the distance between reading and read
frame fail, because they vanish for a self-reading (\(\chi^2\approx34\)).  A single factor on everything read from
the matter frame, \(T_{abc}=(\eta_a-\eta_b)f(c)\), fits the two measurements with \(f(+)=0.944\pm0.008\)
(\(\chi^2=2.9\), one degree of freedom).  In \(\beta\) this is \(0.956\pm0.007\), numerically the known Hopf gap
\(\delta_{\mathrm{CL}}=0.951\) of Remark~\ref{conj:cl-correction}.  The backward analysis yields no new number; it
offers a different reading of a known one.  Only \(f(+)\) is accessible, since every measurement is made from the
matter frame.

\subsection*{v4.8.0-r19 cubes linked by transition maps}
\label{subsec:v480-r19-cubes}

Take one readout cube \(T^{(d)}_{abc}\) per frame \(d\in\{+,0,-\}\) and transition maps
\(\Lambda(d'\!\leftarrow\!d)\) with \(T^{(d')}=\Lambda(d'\!\leftarrow\!d)\,T^{(d)}\).  If the maps are fixed by the
theory, the three cubes carry the information of one cube and the transition law; the family adds no parameter.
What it adds is a consistency question: whether the maps compose to the identity around a closed path of frames.  The
residue, if any, is the holonomy of the family.  It depends on the transition group and the positions of the frames
only, and it is the part that no choice of the reading frame can remove.  If the transition maps commute, the
holonomy is trivial and the family adds nothing.

\subsection*{v4.8.0-r19 the transition group in the two readings}
\label{subsec:v480-r19-transition-group}

\paragraph{Earlier reading.}
Theorem~\ref{thm:symmetric-frame} places matter and tachyon at \(\pm\bsym\) on one axis around the photon and calls
\(\gL\) the Lorentz factor between them.  The transitions are collinear boosts; they commute, and the holonomy is
trivial.

\paragraph{Reading~B.}
By \S\ref{subsec:v480-r06-readings} the kernel is the Riemann--Silberstein vector \(z=E+iB\) under
\(SO(3,\mathbb C)\cong SO^+(1,3)\), with \(s=E^2-B^2+2i\,E\cdot B\).  The working assignment of r07 reads matter as
\(s>0\) (electric type), tachyon as \(s<0\) (magnetic type), photon as \(s=0\) (null field).

\paragraph{Theoremlet v4.8.0-r19.1.}
In Reading~B no Lorentz transformation relates matter to tachyon, and none relates either of them to the photon.
Matter and tachyon with equal \(|s|\) are related by the duality rotation \(z\mapsto e^{i\pi/2}z\), that is
\(E\mapsto-B\), \(B\mapsto E\).  The photon is reached by neither group action.

\paragraph{Proof sketch.}
\(SO(3,\mathbb C)\) preserves \(s\) (Theoremlet r06.3), so it preserves its sign and the null cone.  The duality
rotation multiplies \(s\) by \(e^{2i\alpha}\) and maps \(s>0\) to \(s<0\) at \(\alpha=\pi/2\)
(Theoremlet r07.2); it preserves \(s=0\).  This is the familiar statement that a radiation field is null in every
frame and under every duality rotation.  \(\square\)

\paragraph{Consequence.}
The Lorentz factor \(\gL\) between matter and tachyon in Theorem~\ref{thm:symmetric-frame} has no counterpart in
Reading~B, where the two are related by a duality rotation and not by a boost.  On this point the two readings
disagree about the transition group itself.

\paragraph{Theoremlet v4.8.0-r19.2.}
The duality circle carries a \(\mathbb Z_2\) holonomy around the photon.  Under \(z\mapsto e^{i\alpha}z\) one has
\(s\mapsto e^{2i\alpha}s\); as \(\alpha\) runs from \(0\) to \(\pi\), \(s\) encircles \(s=0\) once and \(z\) returns as
\(-z\).  Paths of \(s\) that do not encircle \(s=0\) have trivial holonomy.  The photon is the branch point of
\(z\) over \(s\).

\paragraph{Proof sketch.}
\(s\) is quadratic in \(z\), so the map from the duality circle of \(z\) to the circle of \(s\) is two to one.
Theoremlet r07.1 states the sign ambiguity of \(z_m\) under \(\alpha\mapsto\alpha+\pi\).  \(\square\)

\subsection*{v4.8.0-r19 the earlier reading as a photon-centred view}
\label{subsec:v480-r19-photon-centred}

\paragraph{Theoremlet v4.8.0-r19.3.}
The real axis of the \(s\)-plane reproduces the sector structure of Theorem~\ref{thm:symmetric-frame}: matter and
tachyon at \(\pm|s|\), mirror-symmetric about the photon at \(s=0\), the mirror \(s\mapsto-s\) being the half-turn of
the duality circle.

\paragraph{Proof sketch.}
Theoremlet r07.1 with \(\alpha=0\) and \(\alpha=\pi/2\).  \(\square\)

\paragraph{Reading, offered and not adopted.}
The photon is fixed by the whole duality circle, so it picks out no reference direction in the \(s\)-plane.  The
earlier reading can then be taken as the view centred on the photon with the real axis chosen as reference, that is,
with the duality angle fixed as a gauge.  On this axis the straight path from matter to tachyon passes through the
photon, the branch point, where the phase is undefined; in Reading~B the path runs around the photon on
\(|s|=\mathrm{const}\).  The photon-centred view therefore cannot see the \(\mathbb Z_2\) holonomy.  It also supplies
something Reading~B lacks: a metric along the axis, in which the mirror distance is read as a boost with factor
\(\gL\).  No correspondence between \(|s|\) and \(\bsym\) is given in the source, and by r19.1 the boost reading has no
counterpart in Reading~B.  Whether the duality angle is a gauge, which would make the earlier reading a gauge-fixed
projection, or physical, which would make it a projection that loses information, is the question left open in r07.

\subsection*{v4.8.0-r19 consequence for the readout factor}
\label{subsec:v480-r19-readout-factor}

The \(\mathbb Z_2\) holonomy can only contribute \(\pm1\); it cannot produce a factor such as \(f(+)=0.944\).  A
continuous holonomy would need closed paths in an \(SO(3,\mathbb C)\)-orbit together with a rule for transporting
frames along them, i.e.\ a connection.  The source supplies neither.  Composing boosts in non-collinear directions
yields a Wigner rotation, and the triole directions are a candidate for such directions, but choosing them is a model
decision.  Any comparison of such a holonomy with \(f(+)\) is a prediction and requires a pre-registration note.

\subsection*{v4.8.0-r19 register entries}
\label{subsec:v480-r19-register}

\begin{center}
\small
\begin{tabular}{@{}>{\raggedright\arraybackslash}p{0.34\textwidth}>{\raggedright\arraybackslash}p{0.58\textwidth}@{}}
\toprule
statement & status after r19\\
\midrule
Thm.~\ref{thm:symmetric-frame} & sector structure reproduced by the real \(s\)-axis of Reading~B (r19.3); \(\gL\) as a Lorentz factor between matter and tachyon has no counterpart in Reading~B (r19.1)\\
duality circle of r07 & \(\mathbb Z_2\) holonomy around the photon, photon as branch point (r19.2)\\
Rem.~\ref{conj:cl-correction}, \(\delta_{\mathrm{CL}}\) & numerically equal to the readout factor \(f(+)\) of the backward analysis; reinterpretation only, not a test (r19)\\
H12 & not separable from H9 by the fixed criterion; no test (r19)\\
\bottomrule
\end{tabular}
\end{center}

\subsection*{v4.8.0-r19 ledger and guard update}
\label{subsec:v480-r19-ledger-guard}

\[
R_{\rm hol}^{v480\text{-}r19}=(0,0,1)
\]
for (transition groups of both readings compared, \(\mathbb Z_2\) holonomy of the duality circle recorded,
continuous holonomy open for lack of a connection).

The layer spends no guard:
\[
R_{\rm glob}^{(v480\text{-}r19,+)}=(0,1,1,1,1),
\qquad
\Gamma_{\rm guard}^{v480\text{-}r19}=(1,1,1,1,1),
\qquad
\beta_4=1/256\ \text{unspent}.
\]
It does not imply exact finite \(C5\) zero closure, exact all-readout terminality,
\(C6/T2\), a global photon-clock/RH theorem, or global all-readout closure.

% ============================================================================
\clearpage
\section*{v4.8.0-r20 Source-Only: The Tachyon Frame, Complex Rapidity and the Holonomy Test H13}
\addcontentsline{toc}{section}{v4.8.0-r20 Source-Only: The Tachyon Frame, Complex Rapidity and the Holonomy Test H13}
\label{sec:v480-r20-tachyon-complex-rapidity}
% ============================================================================

\paragraph{Purpose.}
Layer r19 found no Lorentz transformation between the sectors of Reading~B and read the earlier reading as a
photon-centred view.  This layer records the pre-registered holonomy test H13 and a post-hoc finding attached to it,
and examines the tachyon frame under the holonomy and spinor structure.  The computations are
\texttt{llm/k35\_test\_h13.py} and \texttt{llm/k36\_tachyon\_frame.py}; only H13 is compared with data.  Earlier layers
are not rewritten.

\noindent\textbf{Source-only status.}
This v4.8.0-r20 layer is source-only.  It does not build a PDF, does not assign a new DOI,
and does not authorize a public release.  The v4.8.0 Version DOI
\href{https://doi.org/10.5281/zenodo.23365244}{10.5281/zenodo.23365244} and the Concept DOI
\href{https://doi.org/10.5281/zenodo.19210728}{10.5281/zenodo.19210728} are unchanged by it.

\subsection*{v4.8.0-r20 the holonomy test H13}
\label{subsec:v480-r20-h13}

The note \texttt{llm/Vorab\_Notiz\_H13\_Wigner\_Holonomie.md} fixed, before any computation, one construction: three
matter frames moving with rapidity \(\eta_H=\operatorname{artanh}(1/\sqrt3)\) in the three triole directions, linked by
pure boosts into a closed loop, with \(\cos\delta\) of the resulting Wigner rotation to be compared with the readout
factor \(f(+)=0.944\pm0.008\) of the backward analysis recorded in \S\ref{subsec:v480-r19-records}.

The holonomy is a pure rotation (boost part \(10^{-16}\)) by \(\delta=0.4985\) rad \(=28.56^\circ\), equal to the
angular defect of the equilateral hyperbolic triangle with vertices at the three velocities.  Then
\(\cos\delta=0.878\), at \(z=-7.9\) from \(f(+)\).  By the fixed criterion H13 is refuted for this construction, and the
holonomy route with triole boosts is closed for \(f(+)\).  Since \(f(+)\) is itself a fit, this removes an explanation
of the factor, not the factor.

\paragraph{Post-hoc finding, not registered and not counting.}
Asked after the result: \(\cos^2(\delta/2)=(1+\cos\delta)/2=0.939\) meets \(f(+)\) at \(z=-0.5\).  It is one of four
obvious functions of \(\delta\) (\(\cos\delta\): \(-7.9\); \(\cos(\delta/2)\): \(+3.1\); \(1-\delta/2\pi\): \(-2.8\)), roughly
the rate expected by chance.  Its parts have different standing.  The half angle is forced once the spinor is read
out: by Theoremlet r06.4 (\S\ref{subsec:v480-r06-null-cone}) \(z\) is quadratic in \(\psi\in\mathbb C^2\), so a rotation
of the frames by \(\delta\) turns \(\psi\) by \(\delta/2\).  The square is a Born-type reading of \(|\langle\psi|U\psi\rangle|^2\).
The orientation is a choice: \(|\langle\psi|U\psi\rangle|^2=\cos^2(\delta/2)+\sin^2(\delta/2)\cos^2\theta\), with \(\theta\) the
angle to the rotation axis, which is normal to the triole plane.  The triole photon \(c(1,\omega,\omega^2)\) is a null
field whose direction \(E\times B\) is \((1,1,1)\), the normal of the plane \(z_1+z_2+z_3=0\); its spinor lies on the
axis and acquires only the phase \(\delta/2=14.28^\circ\).  The factor \(\cos^2(\delta/2)\) belongs to a spinor in the plane,
such as the helicity of a matter frame moving in a triole direction.

\subsection*{v4.8.0-r20 the tachyon frame}
\label{subsec:v480-r20-tachyon}

\paragraph{Theoremlet v4.8.0-r20.1.}
The duality rotation \(z\mapsto iz\) commutes with every element of \(SO(3,\mathbb C)\).  Hence every closed path in
the tachyon sector has the same holonomy as its matter image (for the triole loop \(\delta=28.56^\circ\)), and \(z\) and
\(iz\) have the same two principal null directions.  Matter and tachyon with \(z_T=iz_M\) differ only in the overall
phase of the spinor product.

\paragraph{Proof sketch.}
\(iz\) is a scalar multiple of \(z\), and \(SO(3,\mathbb C)\) acts linearly.  Writing \(z\) as a symmetric 2-spinor
\(F_{AB}\) with \(F_{11}=\tfrac12(z_1-iz_2)\), \(F_{22}=\tfrac12(-z_1-iz_2)\), \(F_{12}=-\tfrac12z_3\) (the polarisation of the
map of r06.4), the principal spinors are the roots of \(F_{11}\lambda^2+2F_{12}\lambda+F_{22}\), which do not change when
\(F\) is multiplied by \(i\).  Checked numerically.  \(\square\)

\paragraph{Theoremlet v4.8.0-r20.2.}
For a speed \(v>1\),
\[
\operatorname{artanh}v=\operatorname{artanh}\frac1v+\frac{i\pi}{2}.
\]
The four-velocity at rapidity \(\eta+i\pi/2\) is \(i\,w\) with \(w=(\sinh\eta,\cosh\eta)\) spacelike: the real
four-velocity turns from timelike to spacelike and the factor \(i\) carries the phase.  For \(v=1/\bsym=1.8203\) the real
part is \(\eta_{\rm sym}=0.6174\) and the imaginary part is \(\pi/2\), the duality angle that relates matter and tachyon
in Theoremlet r07.2 (\S\ref{subsec:v480-r07-duality-angle}).

\paragraph{Proof sketch.}
\(\tanh(x+i\pi/2)=\coth x\); \(\cosh(\eta+i\pi/2)=i\sinh\eta\), \(\sinh(\eta+i\pi/2)=i\cosh\eta\).  \(\square\)

\paragraph{The same structure in three places.}
A phase \(\pi/2\) on the representative and a change of sign of the real invariant occur in the duality rotation
\(z_T=iz_M\), \(s\mapsto-s\) (r07.2), in the superluminal rapidity of r20.2, and in the signed tachyon scale
\(\ell_T^2=-\lambda_T^2\) of \S\ref{sec:r82-tachyon-frame-signed-scale}, where \(\ell_T=i\lambda_T\).  The last is an
allowed effective statement of the release scope; r20 records the parallel and adds nothing to it.

\paragraph{Reading, offered and not adopted.}
The earlier reading places the tachyon at \(-\bsym\), the real part of its rapidity; Reading~B keeps the imaginary part
\(\pi/2\) as the duality angle.  The two readings are then related by keeping or dropping one imaginary part, which is
the photon-centred view of \S\ref{subsec:v480-r19-photon-centred} expressed in rapidities.

\subsection*{v4.8.0-r20 cross terms with complex tachyon rapidity}
\label{subsec:v480-r20-cross-terms}

With the tachyon at \(-\eta+i\pi/2\) the cross terms of \S\ref{subsec:v480-r18-cross-terms} become
\begin{center}
\small
\begin{tabular}{@{}lcc@{}}
\toprule
pair & real (r18) & with \(i\pi/2\)\\
\midrule
matter--photon & \(\cosh\eta=1.197\) & \(1.197\)\\
photon--tachyon & \(\cosh\eta=1.197\) & \(-i\sinh\eta=-0.657\,i\)\\
matter--tachyon & \(\cosh2\eta=1.864=\gL\) & \(-i\sinh2\eta=-1.574\,i\)\\
\bottomrule
\end{tabular}
\end{center}
The matter--tachyon term becomes purely imaginary with modulus \(\sinh2\eta=\sqrt{\gL^2-1}\), the proper velocity of the
pair instead of its Lorentz factor; for tachyons \(\cosh\) and \(\sinh\) exchange their roles.  The real factor \(\gL\)
arises only when the imaginary part is dropped.  This is Theoremlet r19.1 in numbers: Reading~B has no real Lorentz factor
between matter and tachyon.

\paragraph{The tribonacci condition.}
The fixed-point condition \(\mu=1/\bsym(\mu)\) of Theorem~\ref{thm:tribonacci} reads, in rapidities,
\(\cosh2\eta=\coth\eta=\tanh(\eta+i\pi/2)\): the Lorentz factor of the opposite term equals the speed of a tachyon with
real rapidity \(\eta\).  This is a reformulation, not a new value.

\subsection*{v4.8.0-r20 register entries}
\label{subsec:v480-r20-register}

\begin{center}
\small
\begin{tabular}{@{}>{\raggedright\arraybackslash}p{0.34\textwidth}>{\raggedright\arraybackslash}p{0.58\textwidth}@{}}
\toprule
statement & status after r20\\
\midrule
H13, readout factor as triole holonomy & refuted for the fixed construction, \(z=-7.9\) (r20)\\
\(\cos^2(\delta/2)\) for \(f(+)\) & post-hoc, not registered, does not count; half angle from r06.4, square and orientation are readings (r20)\\
tachyon sector of Reading~B & same holonomy and same principal spinors as matter (r20.1)\\
Thm.~\ref{thm:symmetric-frame}, tachyon at \(-\bsym\) & real part of the superluminal rapidity \(\eta_{\rm sym}+i\pi/2\); imaginary part equals the duality angle (r20.2)\\
\(\gL\) as matter--tachyon factor & real only with the imaginary part dropped; with it \(-i\sinh2\eta\) (r20)\\
Thm.~\ref{thm:tribonacci} & equivalent to \(\cosh2\eta=\coth\eta\) (r20)\\
\bottomrule
\end{tabular}
\end{center}

\subsection*{v4.8.0-r20 ledger and guard update}
\label{subsec:v480-r20-ledger-guard}

\[
R_{\rm tach}^{v480\text{-}r20}=(0,0,1)
\]
for (H13 decided, tachyon frame related to both readings through the complex rapidity, physical status of the
imaginary part open).

The layer spends no guard:
\[
R_{\rm glob}^{(v480\text{-}r20,+)}=(0,1,1,1,1),
\qquad
\Gamma_{\rm guard}^{v480\text{-}r20}=(1,1,1,1,1),
\qquad
\beta_4=1/256\ \text{unspent}.
\]
It does not imply exact finite \(C5\) zero closure, exact all-readout terminality,
\(C6/T2\), a global photon-clock/RH theorem, or global all-readout closure.
